\chapter{選擇問題}
\label{ch:selection}

% 本章的圖共用的樣式
\tikzset{
    sel-outcell/.style={cell, fill=black!7, text=black!40, draw=black!25},
    sel-sc/.style={draw=black!55, minimum width=0.53cm, minimum height=0.56cm, inner sep=0pt, font=\ttfamily\scriptsize},
    sel-scout/.style={sel-sc, fill=black!7, text=black!40, draw=black!25},
    sel-scmed/.style={sel-sc, draw=optStroke, thick, fill=optFill},
    sel-scpiv/.style={sel-sc, draw=memStroke, thick, fill=memFill},
    sel-tt/.style={draw=black!55, rounded corners=2pt, minimum width=0.62cm, minimum height=0.55cm, inner sep=0pt, font=\ttfamily\small, fill=white},
    sel-ttw/.style={sel-tt, draw=memStroke, thick, fill=memFill},
    sel-tto/.style={sel-tt, draw=optStroke, thick, fill=optFill},
    sel-te/.style={draw=black!50},
    sel-colc/.style={draw=black!45, minimum width=0.8cm, minimum height=0.55cm, inner sep=0pt, font=\ttfamily\small},
    sel-colle/.style={sel-colc, fill=coreFill, draw=coreStroke},
    sel-colge/.style={sel-colc, fill=greenFill, draw=greenStroke},
    sel-colx/.style={sel-colc, fill=memFill, draw=memStroke, very thick},
}
% 分割追蹤的一列：\selrow{y}{說明}{數值}{p}{r}{q}；p..r 以外的格子是灰的，q 是橘色
\newcommand{\selrow}[6]{%
    \node[anchor=east, font=\small] at (0.35, #1) {#2};
    \foreach \v [count=\k] in {#3} {
        \ifnum\k=#6\relax
            \node[hlcell, minimum width=0.78cm, minimum height=0.7cm] at (\k*0.82, #1) {\v};
        \else\ifnum\k<#4\relax
            \node[sel-outcell, minimum width=0.78cm, minimum height=0.7cm] at (\k*0.82, #1) {\v};
        \else\ifnum\k>#5\relax
            \node[sel-outcell, minimum width=0.78cm, minimum height=0.7cm] at (\k*0.82, #1) {\v};
        \else
            \node[cell, minimum width=0.78cm, minimum height=0.7cm] at (\k*0.82, #1) {\v};
        \fi\fi\fi
    }
}

期中考改完，老師請助教公布全班三百零一位同學成績的中位數，順便公布最高分和最低分。
助教的第一個念頭是把成績排序，再拿正中間那一個。 只要一個數，卻把三百多個數全部排
好，是不是做了太多事？ 最高分和最低分一起找，最少要比幾次？ 往下讀之前，先估計一
下。

這一章講\textbf{選擇問題}：從 $n$ 個數裡找出第 $i$ 小的那一個。 先看 $i = 1$、
$i = n$ 這兩個最簡單的情形，算出同時找最大和最小、找第二大最少要比幾次；再講兩個在
$O(n)$ 時間內解決一般情形的演算法。 第一個每次隨機挑一個數把陣列分成兩邊，只往其中
一邊找，期望時間是線性；第二個用「中位數的中位數」挑分割的基準，最壞情況也是線性。
兩個演算法都建立在第~14~章快速排序的分割（partition）上。

這一章用到第~1~章的插入排序和期望值、第~2~章的代換法、第~12~章的指示變數；C 程式
用到《C 程式設計基礎》第~15~章的遞迴和第~19~章的 \texttt{malloc}、\texttt{free}。

\section{第~\texorpdfstring{$i$}{i}~小、最大與最小}
\label{sec:sel-minmax}

$n$ 個數由小到大排好以後，排在第 $i$ 個的數叫做第 $i$ 小的數，也叫第 $i$ 個
\textbf{順序統計量}（order statistic）。 第 $1$ 小是最小值，第 $n$ 小是最大值。
\textbf{中位數}（median）是正中間的那一個：$n$ 是奇數時是第 $(n+1)/2$ 小；$n$ 是偶
數時中間有兩個，第 $n/2$ 小叫\textbf{下中位數}，第 $n/2 + 1$ 小叫\textbf{上中位數}。
這本書說「中位數」時一律指下中位數，也就是第 $\lfloor (n+1)/2 \rfloor$ 小，$n$ 是奇
數時兩種說法一樣。 統計課在 $n$ 是偶數時取兩個中間數的平均，那是另一個定義，在演算
法裡不用，因為平均值不一定是輸入裡的數。

\textbf{選擇問題}（selection problem）：輸入 $n$ 個數和一個整數 $i$，$1 \le i \le n$，
輸出第 $i$ 小的數。 為了敘述方便，這一章的分析都假設 $n$ 個數互不相同。 有重複時，
把每個數換成（數值，原本的位置）這一對，先比數值、數值一樣再比位置，就變成互不相同
了；演算法本身大多不用改，第~\ref{sec:sel-random}~節最後會說明哪裡要小心。

最直接的做法是先排序再取第 $i$ 個，要 $\Theta(n \lg n)$ 的時間。 也可以用第~8~章的
堆積：由下而上建最小堆積花 $O(n)$，再取出 $i$ 次最小值，每次 $O(\lg n)$，共
$O(n + i \lg n)$；$i$ 很小時很快，但找中位數時還是 $\Theta(n \lg n)$。 這一章的目標
是對任何 $i$ 都做到 $O(n)$。

這一章跟第~15~章一樣，只數「兩個輸入元素之間的比較」：演算法只能透過比較得知元素的
大小，比較次數就是成本。 迴圈變數、名次之間的比較不算。

\subsection{找最大值：\texorpdfstring{$n - 1$}{n - 1} 次，而且不能更少}

從左到右掃一次，記住目前看過最大的數，每個新的數和它比一次，共比 $n - 1$ 次。 這已
經是最少了。 把每一次比較想成一場比賽，大的贏、小的輸。 如果演算法結束時，除了它回
答的最大值 $m$ 之外，還有一個數 $y$ 一場都沒輸過，那麼把 $y$ 改成比所有數都大，所有
比較的結果都不會變（$y$ 每一場都是贏的，改大了還是贏），演算法卻還是回答 $m$，答錯
了。 所以除了最大值，其他 $n - 1$ 個數每個都至少要輸一場；一次比較只產生一個輸家，
至少要比 $n - 1$ 次。

\subsection{同時找最大和最小}

分開找要 $2(n - 1)$ 次。 可以省下大約四分之一：把數兩個兩個一組，先讓同一組的兩個
數比一次，小的那個才有資格當最小值，大的那個才有資格當最大值，所以小的只和目前的最
小值比，大的只和目前的最大值比。 每兩個數花 $3$ 次比較，而不是 $4$ 次。

開頭要看 $n$ 的奇偶。 $n$ 是奇數時，第一個數同時當作目前的最小和最大，剩下的偶數個
兩兩一組；$n$ 是偶數時，前兩個數先比一次決定最小和最大，剩下的兩兩一組。

\begin{pseudo}[label=alg:sel-minmax]{同時找最小和最大}
\Function{Min-Max}{$A, n$}
    \If{$n$ 是奇數}
        \State $\mathit{lo} \gets A[1]$；$\mathit{hi} \gets A[1]$；$s \gets 2$
    \Else
        \State 比較 $A[1]$、$A[2]$，小的給 $\mathit{lo}$、大的給 $\mathit{hi}$；$s \gets 3$
    \EndIf
    \For{$j \gets s$ \To $n - 1$，每次加 $2$}
        \State 比較 $A[j]$、$A[j + 1]$，小的給 $u$、大的給 $v$
        \If{$u < \mathit{lo}$}
            \State $\mathit{lo} \gets u$
        \EndIf
        \If{$v > \mathit{hi}$}
            \State $\mathit{hi} \gets v$
        \EndIf
    \EndFor
    \State \Return $(\mathit{lo}, \mathit{hi})$
\EndFunction
\end{pseudo}

\begin{theorem}
\label{thm:sel-minmax}
$n \ge 1$ 時，\proc{Min-Max} 剛好比較 $\lceil 3n/2 \rceil - 2$ 次。
\end{theorem}

\begin{proof}
第 6 行的迴圈每一輪處理 $A[j]$、$A[j+1]$ 兩個數，比較 $3$ 次（第 7、8、10 行）。

$n$ 是奇數：第 3 行不比較，迴圈處理 $A[2..n]$ 的 $(n - 1)/2$ 組，共
$3(n-1)/2$ 次。 而 $\lceil 3n/2 \rceil = (3n + 1)/2$，所以
$\lceil 3n/2 \rceil - 2 = (3n - 3)/2 = 3(n-1)/2$，相等。

$n$ 是偶數：第 5 行比較 $1$ 次，迴圈處理 $A[3..n]$ 的 $(n-2)/2$ 組，共
$1 + 3(n-2)/2 = 3n/2 - 2 = \lceil 3n/2 \rceil - 2$ 次。

正確性：迴圈每一輪開始時，$\mathit{lo}$ 和 $\mathit{hi}$ 是 $A[1..j-1]$ 的最小值和最
大值（初始時 $j = s$，第 2--5 行正好處理了 $A[1..s-1]$）。 一輪裡，$A[j]$ 和
$A[j+1]$ 之中較大的 $v$ 不可能是新的最小值（$u$ 比它小），較小的 $u$ 也不可能是新的
最大值，所以只要拿 $u$ 更新最小值、$v$ 更新最大值，不變量就維持到 $j + 2$。 迴圈結
束時 $j = n + 1$，$\mathit{lo}$、$\mathit{hi}$ 是整個陣列的最小和最大。
\end{proof}

這個次數也是最少的：任何只用比較的演算法，在最壞情況下都至少要比
$\lceil 3n/2 \rceil - 2$ 次。 證明用的是第~15~章證明合併下界時用過的對手論證（adversary argument）：想像
有一個對手在演算法執行時才決定輸入，每次演算法要比較兩個數，對手就挑一個對演算法最
不利、但和之前的回答不矛盾的結果。 只要對手能逼演算法比很多次，就證明了在最壞情況
下一定要比這麼多次，因為對手回答完所有比較之後，一定能寫出一組和這些回答一致的實際
數值，這組數值就是那個「最壞的輸入」。

\begin{whybox}[同時找最大最小的下界]
用前面「每一場比賽」的說法。 演算法結束時，除了最大值之外的 $n - 1$ 個數，每個都必
須輸過至少一場；除了最小值之外的 $n - 1$ 個數，每個都必須贏過至少一場（理由和找最
大值時一樣：一個從沒贏過的數可以被改成比所有數都小）。 把「某個數輸過」「某個數贏
過」各算一條\textbf{消息}，演算法總共要收集至少 $2n - 2$ 條消息。

把每個數分成四種狀態：沒比過（N）、只贏過（W）、只輸過（L）、贏過也輸過（B）。 對
手這樣回答：
\begin{itemize}
\item N 對 N：隨便回答，兩個數分別變成 W 和 L，得到 $2$ 條消息。
\item N 對 W：W 贏；N 對 L：L 輸；N 對 B：N 贏。 都只有原本是 N 的那個數得到 $1$ 條
新消息。
\item W 對 W：其中一個輸，變成 B，$1$ 條。 L 對 L：其中一個贏，變成 B，$1$ 條。
\item W 對 L：W 贏；W 或 L 對 B：W 贏、L 輸；B 對 B：照目前定好的數值回答。 都沒有新
消息。
\end{itemize}
這些回答不會互相矛盾：對手隨時可以把只贏過的數改得比所有數都大、把只輸過的數改得比
所有數都小，而不影響之前的結果，上面每一條規則都只用到這兩種修改。

只有 N 對 N 的比較能一次得到 $2$ 條消息，而每個數只會當一次 N，所以這種比較最多
$\lfloor n/2 \rfloor$ 次，其他比較每次最多 $1$ 條消息。 要收集 $2n - 2$ 條，比較次數
至少
\[
\left\lfloor \frac{n}{2} \right\rfloor + \left(2n - 2 - 2\left\lfloor \frac{n}{2} \right\rfloor\right)
= 2n - 2 - \left\lfloor \frac{n}{2} \right\rfloor
= \left\lceil \frac{3n}{2} \right\rceil - 2 .
\]
最後一步用到 $2n - \lfloor n/2 \rfloor = \lceil 3n/2 \rceil$。
\end{whybox}

\subsection{找第二大：錦標賽}

第二大的數只輸給最大的那一個，所以它一定\textbf{直接}和最大值比過：如果它從沒和最大
值比過，又被判定不是最大，它一定輸給了某個別的數，那個數比它大，它就不是第二大了。
（演算法也必須確認它不是最大，所以它一定至少輸過一場。）

這引出\textbf{錦標賽}（tournament）的做法：

\begin{enumerate}
\item 把 $n$ 個數兩兩配對比賽，贏的晉級；人數是奇數時，最後一個人輪空，直接晉級。
一直比到剩下一個人，就是最大值。 每場比賽淘汰一個人，共 $n - 1$ 場。
\item 第二大一定在「直接輸給冠軍的人」裡面。 冠軍每一輪最多比一場，總共
$\lceil \lg n \rceil$ 輪（每一輪人數變成 $\lceil m/2 \rceil$，從 $n$ 變到 $1$ 要
$\lceil \lg n \rceil$ 輪），所以這些人最多 $\lceil \lg n \rceil$ 個，在他們之中找最大，
最多比 $\lceil \lg n \rceil - 1$ 次。
\end{enumerate}

\begin{center}
\begin{tikzpicture}[x=1cm, y=0.95cm]
\node[sel-tt] (L1) at (0.000,0.00) {47};
\node[sel-tt] (L2) at (1.050,0.00) {65};
\node[sel-tto] (L3) at (2.100,0.00) {88};
\node[sel-ttw] (L4) at (3.150,0.00) {91};
\node[sel-tt] (L5) at (4.200,0.00) {38};
\node[sel-tt] (L6) at (5.250,0.00) {72};
\node[sel-tt] (L7) at (6.300,0.00) {56};
\node[sel-tt] (L8) at (7.350,0.00) {83};
\node[sel-tt] (L9) at (8.400,0.00) {60};
\node[sel-tt] (L10) at (9.450,0.00) {79};
\node[sel-tt] (L11) at (10.500,0.00) {74};
\node[sel-tto] (R1_1) at (0.525,1.00) {65};
\node[sel-ttw] (R1_2) at (2.625,1.00) {91};
\node[sel-tt] (R1_3) at (4.725,1.00) {72};
\node[sel-tt] (R1_4) at (6.825,1.00) {83};
\node[sel-tt] (R1_5) at (8.925,1.00) {79};
\node[sel-tt] (R1_6) at (10.500,1.00) {74};
\node[sel-ttw] (R2_1) at (1.575,2.00) {91};
\node[sel-tto] (R2_2) at (5.775,2.00) {83};
\node[sel-tt] (R2_3) at (9.713,2.00) {79};
\node[sel-ttw] (R3_1) at (3.675,3.00) {91};
\node[sel-tto] (R3_2) at (9.713,3.00) {79};
\node[sel-ttw] (R4_1) at (6.694,4.00) {91};
\draw[sel-te] (R1_1) -- (L1);
\draw[sel-te] (R1_1) -- (L2);
\draw[sel-te] (R1_2) -- (L3);
\draw[sel-te] (R1_2) -- (L4);
\draw[sel-te] (R1_3) -- (L5);
\draw[sel-te] (R1_3) -- (L6);
\draw[sel-te] (R1_4) -- (L7);
\draw[sel-te] (R1_4) -- (L8);
\draw[sel-te] (R1_5) -- (L9);
\draw[sel-te] (R1_5) -- (L10);
\draw[sel-te] (R1_6) -- (L11);
\draw[sel-te] (R2_1) -- (R1_1);
\draw[sel-te] (R2_1) -- (R1_2);
\draw[sel-te] (R2_2) -- (R1_3);
\draw[sel-te] (R2_2) -- (R1_4);
\draw[sel-te] (R2_3) -- (R1_5);
\draw[sel-te] (R2_3) -- (R1_6);
\draw[sel-te] (R3_1) -- (R2_1);
\draw[sel-te] (R3_1) -- (R2_2);
\draw[sel-te] (R3_2) -- (R2_3);
\draw[sel-te] (R4_1) -- (R3_1);
\draw[sel-te] (R4_1) -- (R3_2);
\foreach \t/\lab in {0/參賽, 1/第 1 輪, 2/第 2 輪, 3/第 3 輪, 4/第 4 輪} {
    \node[anchor=east, font=\small] at (-0.5, \t) {\lab};
}
\end{tikzpicture}
\diagcaption{$11$ 個數的錦標賽，每個節點寫的是那場比賽的贏家，只有一個子節點的是輪
空。 橘色是冠軍 $91$ 走過的路，粉紅色是直接輸給它的四個數 $88, 65, 83, 79$。 決賽輸
給冠軍的 $79$ 不是第二大，第一輪就輸給冠軍的 $88$ 才是。}
\label{fig:sel-tournament}
\end{center}

\begin{theorem}
\label{thm:sel-second}
$n \ge 2$ 時，錦標賽在最壞情況下用 $n + \lceil \lg n \rceil - 2$ 次比較找出第二大的數。
\end{theorem}

\begin{proof}
第一步 $n - 1$ 次。 第二步：冠軍在每一輪不是比一場就是輪空，最多
$\lceil \lg n \rceil$ 輪，所以直接輸給它的人 $c \le \lceil \lg n \rceil$ 個；在 $c$ 個
數裡找最大要 $c - 1$ 次。 合計不超過 $n - 1 + \lceil \lg n \rceil - 1$。 正確性就是
本小節開頭的論證：第二大在這 $c$ 個人裡面，而這 $c$ 個人都比冠軍小，其中最大的就是
第二大。
\end{proof}

圖~\ref{fig:sel-tournament} 的 $11$ 個數：第一步比 $5 + 3 + 1 + 1 = 10$ 場，冠軍比了
$4$ 場，四個人裡找最大再比 $3$ 次，共 $13 = 11 + \lceil \lg 11 \rceil - 2$ 次。 冠軍
輪空的次數多時，比較會更少（觀念題第~\ref{sel-q-tour}~題）。

這個次數在最壞情況下也是最少的。 下面只寫證明的骨架。

\begin{whybox}[找第二大的下界（骨架）]
任何找第二大的演算法，除了最大和第二大，其他 $n - 2$ 個數都要輸給某個不是最大值的
數（否則它可能是第二大）；再加上第二大本身要輸給最大值。 設最大值總共直接打敗了
$c$ 個數。 這 $c$ 個數裡，除了第二大以外的 $c - 1$ 個，都還要再輸一場給別人，所以
「輸」的總場數至少是 $(n - 1) + (c - 1)$，每場比較只有一個輸家，比較次數至少
$n + c - 2$。

接著用對手論證逼 $c \ge \lceil \lg n \rceil$。 對手給每個數一個權重，一開始都是 $1$。
兩個都沒輸過的數比賽時，權重大的贏，贏家把輸家的權重加到自己身上，輸家的權重變成
$0$；有一方輸過時，讓沒輸過的那一方贏；兩方都輸過時，回答任何和之前一致的結果。 權
重的總和一直是 $n$，最後全部集中在最大值身上。 最大值每贏一個權重為正的對手，自己
的權重最多變成兩倍（對手的權重不比它大），從 $1$ 變成 $n$ 至少要贏
$\lceil \lg n \rceil$ 場，而且每一場的對手都不一樣（輸過的數權重是 $0$）。 所以
$c \ge \lceil \lg n \rceil$，比較次數至少 $n + \lceil \lg n \rceil - 2$。
\end{whybox}

\begin{examnote}
錦標賽的題目，答案要畫出整棵比賽樹（輪空的人照題目的規則處理，沒有說明時讓最後一個
人輪空），標出冠軍一路打敗的人，再在他們之中找最大。 常見的錯是直接拿決賽的輸家當
第二大，圖~\ref{fig:sel-tournament} 就是反例。 比較次數要分開寫「找冠軍」和「在候選
人中找最大」兩部分。
\end{examnote}


\section{隨機化選擇}
\label{sec:sel-random}

一般的 $i$ 要怎麼找？ 第~14~章的快速排序先挑一個基準（pivot），把比它小的
放左邊、比它大的放右邊，再兩邊各自遞迴排序。 分割完之後，基準所在的位置 $q$ 就是它
排序後的位置：左邊有 $q - p$ 個數比它小。 找第 $i$ 小時只需要知道答案在哪一邊，只往
那一邊遞迴，另一邊丟掉不管。 這個想法是 Hoare 在 1961 年提出的，常叫做
\textbf{快速選擇}（quickselect）。

分割直接用第~14~章第~\ref{sec:so-quick}~節的兩個程序。 \proc{Lomuto-Partition}（虛擬
碼~\ref{alg:so-lomuto}）以 $A[r]$ 為基準，結束時 $A[p..q-1]$ 都 $\le A[q]$，$A[q+1..r]$
都 $> A[q]$，回傳 $q$；過程中基準和其他 $r - p$ 個數各比一次。
\proc{Randomized-Partition}（虛擬碼~\ref{alg:so-rand}）先把 $A[p..r]$ 裡隨機挑的一個數
換到 $A[r]$，再做 Lomuto 分割。 \proc{Random}$(p, r)$ 傳回 $p$ 到 $r$ 之間的一個整數，
每個整數的機率都一樣，每次呼叫彼此獨立。

\begin{pseudo}[label=alg:sel-rselect]{隨機化選擇}
\Function{Randomized-Select}{$A, p, r, i$}
    \If{$p = r$}
        \State \Return $A[p]$ \Comment{只剩一個數，$i$ 一定是 $1$}
    \EndIf
    \State $q \gets$ \Call{Randomized-Partition}{$A, p, r$}
    \State $k \gets q - p + 1$ \Comment{基準在 $A[p..r]$ 裡是第 $k$ 小}
    \If{$i = k$}
        \State \Return $A[q]$
    \ElsIf{$i < k$}
        \State \Return \Call{Randomized-Select}{$A, p, q - 1, i$}
    \Else
        \State \Return \Call{Randomized-Select}{$A, q + 1, r, i - k$}
    \EndIf
\EndFunction
\end{pseudo}

呼叫 \proc{Randomized-Select}$(A, 1, n, i)$ 找整個陣列第 $i$ 小的數。 往右邊找時，左
邊的 $k$ 個數（包括基準）都比答案小，所以在右邊要找的是第 $i - k$ 小，這是第 11 行
最容易漏掉的地方。 往左邊找時名次不變。

\textbf{正確性}對 $n = r - p + 1$ 用強歸納法（觀念題第~\ref{sel-q-correct}~題）。 關
鍵是分割之後 $A[q]$ 前面恰好有 $k - 1$ 個數，而且都比它小（互不相同的假設下），後面
的都比它大，所以 $A[q]$ 是第 $k$ 小；$i < k$ 時答案在左邊、名次還是 $i$，$i > k$ 時
答案在右邊、名次是 $i - k$。 遞迴呼叫的範圍不會是空的：$i < k$ 時左邊有
$k - 1 \ge i \ge 1$ 個數，$i > k$ 時右邊有 $n - k \ge i - k \ge 1$ 個數。

\subsection{手算追蹤：成績的中位數}
\label{sec:sel-trace}

考卷上的選擇題通常不用隨機，而是規定基準「固定取最後一個」，答案才唯一。 這相當於
把第 4 行的 \proc{Randomized-Partition} 換成 \proc{Lomuto-Partition}。 下面用這個版本，找
$13$ 位同學成績的中位數，也就是第 $7$ 小。

\begin{center}
\begin{tikzpicture}
\selrow{0.00}{開始}{59,46,47,80,64,90,62,96,73,58,79,44,54}{1}{13}{0}
\selrow{-0.85}{第 1 次}{46,47,44,54,64,90,62,96,73,58,79,59,80}{1}{13}{4}
\selrow{-1.70}{第 2 次}{46,47,44,54,64,62,73,58,79,59,80,96,90}{5}{13}{11}
\selrow{-2.55}{第 3 次}{46,47,44,54,58,59,73,64,79,62,80,96,90}{5}{10}{6}
\selrow{-3.40}{第 4 次}{46,47,44,54,58,59,62,64,79,73,80,96,90}{7}{10}{7}
\foreach \k in {1,...,13} { \node[idxlbl] at (\k*0.82, -4.05) {\k}; }
\end{tikzpicture}
\diagcaption{基準固定取最後一個，找 $13$ 個成績的第 $7$ 小。 每一列是那一次
\proc{Lomuto-Partition} 之後的陣列：白色格子是這一次分割的範圍 $A[p..r]$，橘色是基準最後的位
置 $q$，灰色是已經丟掉、不會再看的部分。}
\label{fig:sel-score-trace}
\end{center}

\begin{center}
\begin{tabular}{cccccccl}
\toprule
次 & $p$ & $r$ & $i$ & 基準 $A[r]$ & $q$ & $k$ & 下一步 \\
\midrule
1 & 1 & 13 & 7 & 54 & 4 & 4 & $i > k$：往右，$i \gets 7 - 4 = 3$ \\
2 & 5 & 13 & 3 & 80 & 11 & 7 & $i < k$：往左，$i$ 不變 \\
3 & 5 & 10 & 3 & 59 & 6 & 2 & $i > k$：往右，$i \gets 3 - 2 = 1$ \\
4 & 7 & 10 & 1 & 62 & 7 & 1 & $i = k$：回傳 $62$ \\
\bottomrule
\end{tabular}
\end{center}

以第 2 次為例：範圍 $A[5..13]$ 的基準是 $80$，比它小的有 $6$ 個，$80$ 落在第 $11$ 格，是
這個範圍的第 $k = 11 - 5 + 1 = 7$ 小；要找的是第 $3$ 小，所以往左邊 $A[5..10]$ 找，名次
不變。 第 1 次往右時，左邊連同基準的 $4$ 個數都比答案小，名次從 $7$ 變成 $3$。

驗算：排好的成績是 $44, 46, 47, 54, 58, 59, 62, 64, 73, 79, 80, 90, 96$，第 $7$ 個是
$62$。 四次分割分別比了 $12, 8, 5, 3$ 次，共 $28$ 次；用插入排序整個排好，這組資料要
比 $49$ 次（可以用第~1~章插入排序的程式驗證）。

\begin{examnote}
選擇的追蹤題，每次分割寫一列：分割之後的整個陣列、這次的範圍 $p..r$、基準、$q$、
$k = q - p + 1$，以及往哪邊走、新的 $i$。 往右走時 $i$ 要減掉 $k$，往左走時不變，閱卷
時這一欄最容易看出有沒有懂。 題目如果規定別的分割方法（例如第~14~章的 Hoare 分割，
或取第一個當基準），陣列的樣子會不同，照題目規定的分割來做。 名次 $i$ 從 $1$ 開始還
是從 $0$ 開始、中位數取上還是取下，也先看題目的定義。
\end{examnote}

\subsection{最壞情況}

每次都挑到範圍裡最大的數當基準，而要找的是最小值，範圍每次只少一個數，比較次數是
$(n-1) + (n-2) + \cdots + 1 = n(n-1)/2$，最壞情況是 $\Theta(n^2)$。 固定取最後一個當
基準時，已經由小到大排好的陣列找最小值就是這種情況。 隨機挑基準時，沒有哪一個輸入
一定會這樣，要連續好幾次都運氣很差才會發生（觀念題第~\ref{sel-q-worst}~題）。 下面證
明平均起來是線性的。

\subsection{期望時間是線性的}

先說清楚「期望」對什麼取。 這裡的輸入是固定的（任何一組 $n$ 個互不相同的數、任何
一個 $i$），隨機性只來自 \proc{Random} 的結果。 所以下面的結論對\textbf{每一個}輸入
都成立，不需要假設輸入是隨機排列的；這和第~1~章插入排序的平均情況不一樣，那裡是對
所有輸入取平均。

\begin{theorem}
\label{thm:sel-rselect}
對 $n$ 個互不相同的數和任何 $i$，\proc{Randomized-Select} 的期望執行時間是
$\Theta(n)$。
\end{theorem}

\begin{proof}
下界：第一次分割就要看過全部 $n$ 個數，所以至少 $\Omega(n)$。

上界。 令 $T(n)$ 是所有大小為 $n$ 的輸入、所有 $i$ 之中，期望執行時間的最大值；再令
$U(n) = \max_{1 \le m \le n} T(m)$，$U(0) = 0$。 $U$ 是不遞減的，而且
$T(n) \le U(n)$。 取常數 $a$ 使得對 $A[p..r]$ 做一次分割、加上遞迴呼叫以外的其他工作，
都不超過 $a n$（$n = r - p + 1 \ge 2$）；$U(1) = T(1)$ 是常數 $b$。

$n \ge 2$ 時，基準是範圍內每一個數的機率都是 $1/n$，所以它在範圍內是第 $k$ 小的機率
也是 $1/n$，$k = 1, \ldots, n$。 已知基準是第 $k$ 小時，接下來最多遞迴一次，範圍大
小是 $k - 1$ 或 $n - k$，之後用到的隨機數是新的，所以這部分的期望時間不超過
$U(\max(k - 1, n - k))$。 對 $k$ 取平均：
\[
T(n) \le a n + \frac{1}{n} \sum_{k=1}^{n} U(\max(k - 1, n - k)).
\]
$k$ 從 $1$ 跑到 $n$ 時，$\max(k - 1, n - k)$ 取到 $\lfloor n/2 \rfloor, \ldots, n-1$
裡的值，每個值最多出現兩次，所以
\[
T(n) \le a n + \frac{2}{n} \sum_{m=\lfloor n/2 \rfloor}^{n-1} U(m).
\]

用代換法（第~2~章第~\ref{sec:rec-subst}~節）證明：取 $c = \max(4a, b)$，所有 $n \ge 1$ 都有
$U(n) \le cn$。 $n = 1$ 時 $U(1) = b \le c$。 設 $n \ge 2$，比 $n$ 小的都成立。 先估
計和：
\[
\sum_{m=\lfloor n/2 \rfloor}^{n-1} m
= \frac{n(n-1)}{2} - \frac{(\lfloor n/2 \rfloor - 1)\lfloor n/2 \rfloor}{2}
\le \frac{n(n-1)}{2} - \frac{(n-3)(n-1)}{8}
= \frac{3(n^2 - 1)}{8} < \frac{3n^2}{8},
\]
中間用到 $\lfloor n/2 \rfloor \ge (n-1)/2 \ge 1/2$，以及 $t(t-1)$ 在 $t \ge 1/2$ 時隨 $t$ 遞增。
代進去：
\[
T(n) \le a n + \frac{2c}{n} \cdot \frac{3n^2}{8} = \frac{3}{4} c n + a n \le c n,
\]
最後一步用到 $c \ge 4a$。 再加上 $U(n - 1) \le c(n-1) \le cn$，得到
$U(n) = \max(U(n-1), T(n)) \le cn$。 所以期望時間 $T(n) \le U(n) = O(n)$。
\end{proof}

證明裡要的是 $U(n) \le cn$ 這個\textbf{同一個} $c$，推出來的是
$\frac{3}{4}cn + an$，多出來的 $an$ 剛好被省下的 $\frac{1}{4}cn$ 吸收。 這就是第~2~章
提醒過的「歸納步驟要推出同一個常數」。 直覺上，基準落在中間一半的機率是 $1/2$，落
在那裡時範圍至少縮成 $3/4$；平均兩次分割就有一次「好的」，總工作量是一個公比小於
$1$ 的等比級數。

\subsection{比較次數算到常數}

上面的證明只說是 $O(n)$。 用第~12~章第~\ref{sec:hash-chain-analysis}~節的指示變數，可以把期望比較次數直接寫
成一個和，再算出常數。

\begin{theorem}
\label{thm:sel-rselect-cmp}
把 $n$ 個互不相同的數由小到大記作 $z_1 < z_2 < \cdots < z_n$。
\proc{Randomized-Select} 找第 $i$ 小時，期望比較次數是
\[
E_n(i) = \sum_{1 \le a < b \le n} \frac{2}{\max(b, i) - \min(a, i) + 1},
\]
而且 $E_n(i) < 4n$。
\end{theorem}

\begin{whybox}[期望比較次數]
比較只發生在分割時，基準和範圍裡的其他數各比一次；當過基準的數之後就不在任何範圍
裡了，所以每一對 $z_a, z_b$ 最多比一次。 令 $X_{ab}$ 是「$z_a$ 和 $z_b$ 比過」的指
示變數，總比較次數是 $\sum_{a<b} X_{ab}$，期望值是 $\sum_{a<b} \Pr[X_{ab} = 1]$。

固定 $a < b$，令 $Z$ 是名次從 $\min(a, i)$ 到 $\max(b, i)$ 的那些數，共
$\max(b, i) - \min(a, i) + 1$ 個；$Z$ 同時包含 $z_a$、$z_b$ 和答案 $z_i$。
\begin{itemize}
\item 還沒有 $Z$ 裡的數被選成基準時，$Z$ 整個都在目前的範圍裡：一個 $Z$ 以外的基準
不是比 $Z$ 全部小，就是比 $Z$ 全部大，而演算法留下的是含有 $z_i$ 的那一邊，也就是
含有整個 $Z$ 的那一邊。
\item 所以第一個落在 $Z$ 裡的基準，在 $Z$ 的每個數上機率都一樣，是 $z_a$ 或 $z_b$ 的
機率是 $2/|Z|$。 是的話，$z_a$、$z_b$ 這時都在範圍裡，比了一次。
\item 如果它是 $Z$ 裡的其他數 $z_c$：$a < c < b$ 時，$z_a$ 和 $z_b$ 被分到兩邊（或
$c = i$，演算法直接結束），以後不會再比；$c$ 在 $a$、$b$ 的範圍外時，$c$ 夾在 $i$ 和
這一對之間（或 $c = i$），留下的那一邊含有 $z_i$ 而不含 $z_a$、$z_b$，也不會再比。
\end{itemize}
所以 $\Pr[X_{ab} = 1] = 2/(\max(b,i) - \min(a,i) + 1)$，加起來就是 $E_n(i)$。

上界分三種 $(a, b)$ 估計：
\begin{itemize}
\item $a \le i \le b$：分母是 $b - a + 1 = L$。 長度為 $L$、包含 $i$ 的區間最多 $L$ 個，
每個貢獻 $2/L$，所以長度 $L$ 的總貢獻不超過 $2$；$L = 2, \ldots, n$，合計不超過
$2(n - 1)$。
\item $i < a < b$：分母是 $b - i + 1$。 固定 $b$，$a$ 有 $b - i - 1$ 種，貢獻
$2(b - i - 1)/(b - i + 1) < 2$；$b$ 最多 $n - i$ 種，合計小於 $2(n - i)$。
\item $a < b < i$：同理，合計小於 $2(i - 1)$。
\end{itemize}
三項加起來小於 $2(n-1) + 2(n-i) + 2(i-1) < 4n$。
\end{whybox}

找最小值（$i = 1$）時，這個和可以算出來是 $2n - 2H_n$（觀念題第~\ref{sel-q-min}~題），
比直接掃一次的 $n - 1$ 多了將近一倍：隨機選擇不是用來找最大最小的。 找中位數時和最
大，$n$ 很大時大約是 $3.39n$。 下面的程式實際跑一次，和公式對照。

\subsection{C 程式}

C 標準函式庫的 \texttt{rand} 在不同的系統上產生的數列不一樣，書上印的輸出就對不起來。
所以程式自己寫一個\textbf{線性同餘產生器}（linear congruential generator，LCG）：
$s \gets s \cdot M + C \bmod 2^{64}$，$M$、$C$ 是固定的常數，種子 $s$ 一開始固定，
每次執行產生的數列都一樣。 LCG 的低位元週期很短，所以取高位元當結果。
\texttt{random\_between(p, r)} 用取餘數把結果縮到 $p..r$，會有一點點不均勻（《C 程式
設計基礎》第~18~章講的模偏差），範圍遠小於 $2^{31}$ 時可以忽略。

\begin{cbookrecap}[第~13~章 13.1、13.5 節，第~18~章 18.8 節]
有號整數溢位是未定義行為；\texttt{unsigned} 型別的運算保證是模 $2^n$，所以 LCG 的乘
法和加法要用 \texttt{unsigned long long}，溢位就是我們要的「模 $2^{64}$」。 寫在函式
外面的 \texttt{static} 變數只在這個檔案裡看得到，而且整個程式執行期間都存在，適合存
產生器的狀態。 \texttt{rand()} 的數列由實作決定，換一台機器就可能不同。
\end{cbookrecap}

程式先用隨機分割找第~\ref{sec:sel-trace}~節那 $13$ 個成績的中位數；再拿 $1001$ 個不同的數，同一
個輸入跑 $2000$ 次，算三個名次的平均比較次數，和定理~\ref{thm:sel-rselect-cmp} 的公式
比較。 \texttt{ncmp} 只在分割的迴圈裡加一，正好是元素之間的比較次數。

\begin{example}[label=ex:sel-rselect]{c}{(rselect.c) 隨機化選擇與期望比較次數}
#include <stdio.h>
#include <stdlib.h>

// 自己寫的線性同餘產生器（LCG），種子固定，每次執行的結果都一樣。
static unsigned long long state = 20261119ULL;

int random_between(int p, int r)            // p..r 之間的一個整數
{
    state = state * 6364136223846793005ULL + 1442695040888963407ULL;
    // 低位元的週期短，取高位元
    unsigned x = (unsigned)(state >> 33);
    return p + (int)(x % (unsigned)(r - p + 1));
}

long long ncmp;                             // 元素之間的比較次數

void swap(int *x, int *y)
{
    int t = *x;
    *x = *y;
    *y = t;
}

// 隨機挑一個數換到 a[r] 當基準，做 Lomuto 分割，回傳基準最後的位置。
int randomized_partition(int *a, int p, int r)
{
    swap(&a[r], &a[random_between(p, r)]);
    int x = a[r], i = p - 1;
    for (int j = p; j < r; j++) {
        ncmp++;
        if (a[j] <= x)
            swap(&a[++i], &a[j]);
    }
    swap(&a[i + 1], &a[r]);
    return i + 1;
}

// 回傳 a[p..r] 裡第 i 小的數，1 <= i <= r - p + 1。
int randomized_select(int *a, int p, int r, int i)
{
    if (p == r)
        return a[p];
    int q = randomized_partition(a, p, r);
    int k = q - p + 1;
    if (i == k)
        return a[q];
    else if (i < k)
        return randomized_select(a, p, q - 1, i);
    else
        return randomized_select(a, q + 1, r, i - k);
}

// 定理的公式：n 個不同的數找第 k 小時的期望比較次數。
double expected_cmp(int n, int k)
{
    double s = 0;
    for (int a = 1; a <= n; a++)
        for (int b = a + 1; b <= n; b++)
            s += 2.0 / ((b > k ? b : k) - (a < k ? a : k) + 1);
    return s;
}

int main(void)
{
    int score[] = {59, 46, 47, 80, 64, 90, 62, 96, 73, 58, 79,
                   44, 54};
    int med = randomized_select(score, 0, 12, 7);
    printf("median = %d, comparisons = %lld\n", med, ncmp);

    // 1001 個不同的數，同一個輸入跑 2000 次，看平均比較次數。
    int N = 1001, T = 2000, ks[] = {1, 251, 501};
    int *base = malloc(N * sizeof *base);
    int *b = malloc(N * sizeof *b);
    if (base == NULL || b == NULL)
        return 1;
    for (int j = 0; j < N; j++)
        base[j] = j + 1;
    // 洗牌，得到 1..N 的一個排列
    for (int j = N - 1; j > 0; j--)
        swap(&base[j], &base[random_between(0, j)]);
    printf("k average formula 4n\n");
    for (int q = 0; q < 3; q++) {
        long long total = 0;
        for (int t = 0; t < T; t++) {
            for (int j = 0; j < N; j++)
                b[j] = base[j];
            ncmp = 0;
            if (randomized_select(b, 0, N - 1, ks[q]) != ks[q])
                printf("WRONG\n");
            total += ncmp;
        }
        printf("%d %.1f %.1f %d\n", ks[q], (double)total / T,
               expected_cmp(N, ks[q]), 4 * N);
    }
    free(base);
    free(b);
    return 0;
}
\end{example}
\iosample{}{median = 62, comparisons = 22 \\ k average formula 4n \\ 1 1980.5 1987.0 4004 \\ 251 3055.8 3079.8 4004 \\ 501 3349.0 3339.7 4004}

同一組成績，隨機挑基準這一次比了 $22$ 次，比固定取最後一個的 $28$ 次少；換一個種子
結果會不同，但答案一定是 $62$。 實驗的部分，$2000$ 次的平均和公式差不到百分之一，三
個名次都小於 $4n$；名次越靠中間，比較越多，中位數大約是 $3.34n$。

\subsection{有重複的數時}

前面的分析假設數互不相同。 \proc{Randomized-Select} 遇到重複的數仍然回傳正確的答
案（分割後左邊都 $\le$ 基準、右邊都 $>$ 基準，論證照樣成立），但時間可能變差。 最極
端的是 $n$ 個數全部一樣：Lomuto 分割每次都把基準放到範圍的最後一格，找中位數要
$\Theta(n^2)$，隨機挑基準也沒有用。 解法是改成第~14~章的三路分割，把範圍分成小於、
等於、大於基準三段，$i$ 落在等於那一段就直接回傳（觀念題第~\ref{sel-q-dup}~題）。

\section{中位數的中位數：最壞情況線性}
\label{sec:sel-mom}

隨機化選擇的期望是線性，最壞還是 $\Theta(n^2)$。 1973 年 Blum、Floyd、Pratt、Rivest、
Tarjan 給了一個最壞情況也是 $O(n)$ 的演算法。 前一節的論證已經指出關鍵：只要基準保
證不太靠兩端，每次都能丟掉固定比例的數，總工作量就是公比小於 $1$ 的等比級數。 問題
是怎麼不靠運氣、又只花線性時間，就挑到這樣的基準。

做法如下（$n \le 5$ 時直接排序）：

\begin{enumerate}
\item 把 $A[p..r]$ 從左到右每 $5$ 個分成一組，最後一組可能不滿 $5$ 個，共
$g = \lceil n/5 \rceil$ 組。
\item 每組用插入排序排好，取出每組的中位數（不滿的那組取下中位數），放進新陣列
$M[1..g]$。
\item \textbf{遞迴}找出 $M$ 的中位數 $x$，這就是「中位數的中位數」。
\item 用 $x$ 當基準分割 $A[p..r]$，接下來和 \proc{Randomized-Select} 一樣，只往一邊遞
迴。
\end{enumerate}

\begin{pseudo}[label=alg:sel-mom, float=!tp]{中位數的中位數}
\Function{Select}{$A, p, r, i$}
    \State $n \gets r - p + 1$
    \If{$n \le 5$}
        \State 用插入排序把 $A[p..r]$ 排好
        \State \Return $A[p + i - 1]$
    \EndIf
    \State $g \gets \lceil n/5 \rceil$
    \State 建立陣列 $M[1..g]$
    \For{$j \gets 1$ \To $g$}
        \State $\ell \gets p + 5(j - 1)$
        \State $h \gets \min(\ell + 4, r)$
        \State 用插入排序把 $A[\ell..h]$ 排好
        \State $M[j] \gets A[\lfloor (\ell + h)/2 \rfloor]$ \Comment{這一組的（下）中位數}
    \EndFor
    \State $x \gets$ \Call{Select}{$M, 1, g, \lceil g/2 \rceil$}
    \State $q \gets$ \Call{Partition-Around}{$A, p, r, x$}
    \State $k \gets q - p + 1$
    \If{$i = k$}
        \State \Return $A[q]$
    \ElsIf{$i < k$}
        \State \Return \Call{Select}{$A, p, q - 1, i$}
    \Else
        \State \Return \Call{Select}{$A, q + 1, r, i - k$}
    \EndIf
\EndFunction
\Statex
\Function{Partition-Around}{$A, p, r, x$}
    \State 找出 $A[p..r]$ 裡等於 $x$ 的位置 $j$
    \State 交換 $A[j]$ 和 $A[r]$
    \State \Return \Call{Lomuto-Partition}{$A, p, r$}
\EndFunction
\end{pseudo}

一組 $h - \ell + 1$ 個數排好之後，下中位數在第 $\lceil (h - \ell + 1)/2 \rceil$ 個，
位置是 $\ell + \lfloor (h - \ell)/2 \rfloor = \lfloor (\ell + h)/2 \rfloor$，這就是第
12 行。 第 13 行對 $M$ 遞迴，$M$ 只有 $\lceil n/5 \rceil$ 個數；第 14 行花 $O(n)$
找到 $x$ 的位置，再交給 Lomuto 分割。 第 4、11 行排的都是不超過 $5$ 個數，每次比較
不超過 $10$ 次。

\subsection{\texorpdfstring{$x$}{x} 保證不靠兩端}

把每一組排好、立起來當成一欄，小的在下、大的在上，再把各欄依照中位數由小到大排好，
如圖~\ref{fig:sel-mom-cols}。 中間那一列是各組的中位數，$x$ 在正中間。 左下方藍色的
數：它們所在的那一欄的中位數 $\le x$，而它們又 $\le$ 那一欄的中位數，所以都 $\le x$。
同理，右上方綠色的數都 $\ge x$。 不管其他數怎麼排，這兩塊都保證存在。

\begin{figure}[tbp]
\centering
\begin{tikzpicture}
\node[sel-colle] at (0.00,0.00) {16};
\node[sel-colle] at (0.00,0.62) {27};
\node[sel-colle] at (0.00,1.24) {31};
\node[sel-colc] at (0.00,1.86) {70};
\node[sel-colc] at (0.00,2.48) {87};
\node[idxlbl] at (0.00,-0.55) {第 4 組};
\node[sel-colle] at (0.95,0.00) {13};
\node[sel-colle] at (0.95,0.62) {30};
\node[sel-colle] at (0.95,1.24) {35};
\node[sel-colc] at (0.95,1.86) {66};
\node[sel-colc] at (0.95,2.48) {85};
\node[idxlbl] at (0.95,-0.55) {第 2 組};
\node[sel-colle] at (1.90,0.00) {15};
\node[sel-colle] at (1.90,0.62) {33};
\node[sel-colx] at (1.90,1.24) {38};
\node[sel-colge] at (1.90,1.86) {88};
\node[sel-colge] at (1.90,2.48) {98};
\node[idxlbl] at (1.90,-0.55) {第 3 組};
\node[sel-colc] at (2.85,0.00) {24};
\node[sel-colc] at (2.85,0.62) {34};
\node[sel-colge] at (2.85,1.24) {50};
\node[sel-colge] at (2.85,1.86) {71};
\node[sel-colge] at (2.85,2.48) {86};
\node[idxlbl] at (2.85,-0.55) {第 5 組};
\node[sel-colc] at (3.80,0.00) {44};
\node[sel-colc] at (3.80,0.62) {48};
\node[sel-colge] at (3.80,1.24) {80};
\node[sel-colge] at (3.80,1.86) {89};
\node[sel-colge] at (3.80,2.48) {91};
\node[idxlbl] at (3.80,-0.55) {第 1 組};
\draw[diagarrow] (-0.85,-0.2) -- node[left, font=\scriptsize, align=center] {由小\\到大} (-0.85,2.7);
\draw[diagarrow] (-0.4,3.25) -- node[above, font=\scriptsize] {中位數由小到大} (4.2,3.25);
\end{tikzpicture}
\diagcaption{第~\ref{sec:sel-mom-trace}~節第一層的 $25$ 個數。 每一欄是排好的一組，欄
的順序依中位數排（演算法本身不搬欄，這只是為了看清楚）。 $x = 38$（橘色）。 藍色的
$8$ 個數一定 $< x$，綠色的 $8$ 個數一定 $> x$；白色的數可能在 $x$ 的任何一邊。}
\label{fig:sel-mom-cols}
\end{figure}

\begin{lemma}
\label{lem:sel-mom-split}
設 $n \ge 6$，$n$ 個數互不相同。 \proc{Select} 第 14 行分割之後，$x$ 的左邊和右邊各有不超
過 $7n/10 + 2$ 個數。 所以第 19 行或第 21 行遞迴的範圍不超過 $7n/10 + 2$。
\end{lemma}

\begin{proof}
$g = \lceil n/5 \rceil$，$x$ 是 $g$ 個組中位數裡第 $\lceil g/2 \rceil$ 小的。 只有最後
一組可能不滿。

數 $\ge x$ 的數。 中位數 $\ge x$ 的組有 $g - \lceil g/2 \rceil + 1 = \lfloor g/2 \rfloor + 1$
組（包括 $x$ 自己那組）。 滿的一組裡，中位數和比它大的兩個數，共 $3$ 個都 $\ge x$。
不滿的那組有 $s$ 個數（$1 \le s \le 4$），$\ge$ 下中位數的有 $s - \lceil s/2 \rceil + 1$
個，至少 $1$ 個，比 $3$ 最多少 $2$。 所以 $\ge x$ 的數至少
$3(\lfloor g/2 \rfloor + 1) - 2 = 3\lfloor g/2 \rfloor + 1$ 個，$x$ 左邊（$< x$）的數
最多
\[
n - 3\left\lfloor \frac{g}{2} \right\rfloor - 1
\le n - 3 \cdot \frac{g - 1}{2} - 1
\le n - \frac{3}{2}\left(\frac{n}{5} - 1\right) - 1
= \frac{7n}{10} + \frac{1}{2}.
\]

數 $\le x$ 的數。 中位數 $\le x$ 的組有 $\lceil g/2 \rceil$ 組，滿的組各貢獻 $3$ 個，
不滿的組至少 $\lceil s/2 \rceil \ge 1$ 個，一樣最多少 $2$。 所以 $\le x$ 的至少
$3\lceil g/2 \rceil - 2$ 個，$x$ 右邊的數最多
\[
n - 3\left\lceil \frac{g}{2} \right\rceil + 2 \le n - \frac{3n}{10} + 2 = \frac{7n}{10} + 2,
\]
用到 $\lceil g/2 \rceil \ge g/2 \ge n/10$。 兩邊都不超過 $7n/10 + 2$。
\end{proof}

圖~\ref{fig:sel-mom-cols} 的例子：$n = 25$，引理保證每一邊不超過 $19$ 個；實際上藍色
和綠色各 $8$ 個，連同 $x$ 自己，每一邊都保證有 $9$ 個數不在那一邊，所以每一邊不超過
$25 - 9 = 16$ 個。

\subsection{最壞情況是線性的}

\begin{theorem}
\label{thm:sel-mom}
對 $n$ 個互不相同的數，\proc{Select} 的最壞執行時間是 $O(n)$。
\end{theorem}

\begin{proof}
令 $T(n)$ 是大小為 $n$ 的所有輸入、所有 $i$ 之中最長的執行時間。 第 13 行的遞迴大小
是 $\lceil n/5 \rceil$；分組、排各組、找 $x$ 的位置、分割，合計不超過 $an$（$a$ 是常
數）；由引理，第 19、21 行最多一次遞迴，大小不超過 $7n/10 + 2$。 所以 $n \ge 6$ 時
\[
T(n) \le T(\lceil n/5 \rceil) + T(m) + an, \qquad m \le \frac{7n}{10} + 2,
\]
沒有遞迴（$i = k$）時 $T(m)$ 那一項不出現。 兩個遞迴的大小都比 $n$ 小（引理證明裡的
兩個上界 $n - 3\lfloor g/2 \rfloor - 1$、$n - 3\lceil g/2 \rceil + 2$ 都小於 $n$），所以
演算法一定會停，每個 $T(n)$ 都是有限的。

用代換法證明存在常數 $c$，所有 $n \ge 1$ 都有 $T(n) \le cn$。 $n \ge 60$ 時對兩個遞
迴用歸納假設：
\[
T(n) \le c\left(\frac{n}{5} + 1\right) + c\left(\frac{7n}{10} + 2\right) + an
= \frac{9}{10}cn + 3c + an.
\]
要這個 $\le cn$，等於要 $\frac{c}{10}(n - 30) \ge an$。 $n \ge 60$ 時
$n - 30 \ge n/2$，左邊至少 $\frac{c}{20} n$，只要 $c \ge 20a$ 就成立。

$1 \le n < 60$ 的有限多個 $n$，$T(n)$ 都不超過某個常數 $B$，取 $c \ge B$ 就有
$T(n) \le B \le cn$。 所以 $c = \max(20a, B)$ 時，對所有 $n$ 都成立，$T(n) = O(n)$。
\end{proof}

證明能過，是因為兩個遞迴的大小加起來是 $\frac{n}{5} + \frac{7n}{10} = \frac{9n}{10}$：
用第~2~章第~\ref{sec:rec-tree}~節的遞迴樹來看，每往下一層總工作量最多是上一層的 $9/10$，各層加起來是
公比 $9/10$ 的等比級數。 門檻 $60$ 只是讓常數好算。

\subsection{為什麼是 5 個一組}

\textbf{3 個一組}：$\ge x$ 的數大約只有 $2 \cdot \frac{g}{2} = \frac{n}{3}$ 個，往下
遞迴的那一邊可能有 $\frac{2n}{3}$ 個，遞迴式變成
\[
T(n) \le T(n/3) + T(2n/3) + an.
\]
$\frac{1}{3} + \frac{2}{3} = 1$，每一層的工作量都不會變少。 代換法試 $T(n) \le cn$
會得到 $cn + an$，永遠多一個 $an$，常數再大也證不過；這個遞迴式的解是
$\Theta(n \lg n)$（第~2~章第~\ref{sec:rec-tree}~節「不平衡的樹」，把 $1/4 : 3/4$ 換成 $1/3 : 2/3$）。
所以用這個分析，$3$ 個一組只能保證 $O(n \lg n)$。

\textbf{7 個一組}：每組貢獻 $4$ 個，$\ge x$ 的數大約 $4 \cdot \frac{g}{2} = \frac{2n}{7}$
個，往下遞迴的一邊大約 $\frac{5n}{7}$，遞迴式是 $T(n/7) + T(5n/7) + an$，
$\frac{1}{7} + \frac{5}{7} = \frac{6}{7} < 1$，一樣是線性。 組越大，找 $x$ 的遞迴越
小，但每組排序的成本越高；$5$ 是能讓兩個比例加起來小於 $1$ 的最小奇數。

\subsection{手算追蹤：25 個數}
\label{sec:sel-mom-trace}

$A = \langle 48, 89, 80, 91, 44, 13, 35, 30, 85, 66, 98, 33, 38, 88, 15, 70, 87, 31,
16, 27, 24, 50, 86, 71, 34 \rangle$，找中位數，也就是第 $13$ 小。 下圖每一列是一個階
段之後的陣列，粗線是組的分界，粉紅色是各組的中位數，橘色是這一層的 $x$（或最後的答
案），灰色是已經丟掉的部分。

\begin{center}
\begin{tikzpicture}
\node[anchor=east, font=\small] at (0.25,0.00) {原始};
\node[sel-sc] at (0.55,0.00) {48};
\node[sel-sc] at (1.10,0.00) {89};
\node[sel-sc] at (1.65,0.00) {80};
\node[sel-sc] at (2.20,0.00) {91};
\node[sel-sc] at (2.75,0.00) {44};
\node[sel-sc] at (3.30,0.00) {13};
\node[sel-sc] at (3.85,0.00) {35};
\node[sel-sc] at (4.40,0.00) {30};
\node[sel-sc] at (4.95,0.00) {85};
\node[sel-sc] at (5.50,0.00) {66};
\node[sel-sc] at (6.05,0.00) {98};
\node[sel-sc] at (6.60,0.00) {33};
\node[sel-sc] at (7.15,0.00) {38};
\node[sel-sc] at (7.70,0.00) {88};
\node[sel-sc] at (8.25,0.00) {15};
\node[sel-sc] at (8.80,0.00) {70};
\node[sel-sc] at (9.35,0.00) {87};
\node[sel-sc] at (9.90,0.00) {31};
\node[sel-sc] at (10.45,0.00) {16};
\node[sel-sc] at (11.00,0.00) {27};
\node[sel-sc] at (11.55,0.00) {24};
\node[sel-sc] at (12.10,0.00) {50};
\node[sel-sc] at (12.65,0.00) {86};
\node[sel-sc] at (13.20,0.00) {71};
\node[sel-sc] at (13.75,0.00) {34};
\draw[very thick, black!80] (3.025,-0.36) -- (3.025,0.36);
\draw[very thick, black!80] (5.775,-0.36) -- (5.775,0.36);
\draw[very thick, black!80] (8.525,-0.36) -- (8.525,0.36);
\draw[very thick, black!80] (11.275,-0.36) -- (11.275,0.36);
\node[anchor=east, font=\small] at (0.25,-0.95) {各組排好};
\node[sel-sc] at (0.55,-0.95) {44};
\node[sel-sc] at (1.10,-0.95) {48};
\node[sel-scmed] at (1.65,-0.95) {80};
\node[sel-sc] at (2.20,-0.95) {89};
\node[sel-sc] at (2.75,-0.95) {91};
\node[sel-sc] at (3.30,-0.95) {13};
\node[sel-sc] at (3.85,-0.95) {30};
\node[sel-scmed] at (4.40,-0.95) {35};
\node[sel-sc] at (4.95,-0.95) {66};
\node[sel-sc] at (5.50,-0.95) {85};
\node[sel-sc] at (6.05,-0.95) {15};
\node[sel-sc] at (6.60,-0.95) {33};
\node[sel-scpiv] at (7.15,-0.95) {38};
\node[sel-sc] at (7.70,-0.95) {88};
\node[sel-sc] at (8.25,-0.95) {98};
\node[sel-sc] at (8.80,-0.95) {16};
\node[sel-sc] at (9.35,-0.95) {27};
\node[sel-scmed] at (9.90,-0.95) {31};
\node[sel-sc] at (10.45,-0.95) {70};
\node[sel-sc] at (11.00,-0.95) {87};
\node[sel-sc] at (11.55,-0.95) {24};
\node[sel-sc] at (12.10,-0.95) {34};
\node[sel-scmed] at (12.65,-0.95) {50};
\node[sel-sc] at (13.20,-0.95) {71};
\node[sel-sc] at (13.75,-0.95) {86};
\draw[very thick, black!80] (3.025,-1.31) -- (3.025,-0.59);
\draw[very thick, black!80] (5.775,-1.31) -- (5.775,-0.59);
\draw[very thick, black!80] (8.525,-1.31) -- (8.525,-0.59);
\draw[very thick, black!80] (11.275,-1.31) -- (11.275,-0.59);
\node[anchor=east, font=\small] at (0.25,-1.90) {對 38 分割};
\node[sel-sc] at (0.55,-1.90) {13};
\node[sel-sc] at (1.10,-1.90) {30};
\node[sel-sc] at (1.65,-1.90) {35};
\node[sel-sc] at (2.20,-1.90) {15};
\node[sel-sc] at (2.75,-1.90) {33};
\node[sel-sc] at (3.30,-1.90) {16};
\node[sel-sc] at (3.85,-1.90) {27};
\node[sel-sc] at (4.40,-1.90) {31};
\node[sel-sc] at (4.95,-1.90) {24};
\node[sel-sc] at (5.50,-1.90) {34};
\node[sel-scpiv] at (6.05,-1.90) {38};
\node[sel-sc] at (6.60,-1.90) {91};
\node[sel-sc] at (7.15,-1.90) {86};
\node[sel-sc] at (7.70,-1.90) {88};
\node[sel-sc] at (8.25,-1.90) {98};
\node[sel-sc] at (8.80,-1.90) {44};
\node[sel-sc] at (9.35,-1.90) {48};
\node[sel-sc] at (9.90,-1.90) {80};
\node[sel-sc] at (10.45,-1.90) {70};
\node[sel-sc] at (11.00,-1.90) {87};
\node[sel-sc] at (11.55,-1.90) {66};
\node[sel-sc] at (12.10,-1.90) {85};
\node[sel-sc] at (12.65,-1.90) {50};
\node[sel-sc] at (13.20,-1.90) {71};
\node[sel-sc] at (13.75,-1.90) {89};
\node[anchor=east, font=\small] at (0.25,-2.85) {各組排好};
\node[sel-scout] at (0.55,-2.85) {13};
\node[sel-scout] at (1.10,-2.85) {30};
\node[sel-scout] at (1.65,-2.85) {35};
\node[sel-scout] at (2.20,-2.85) {15};
\node[sel-scout] at (2.75,-2.85) {33};
\node[sel-scout] at (3.30,-2.85) {16};
\node[sel-scout] at (3.85,-2.85) {27};
\node[sel-scout] at (4.40,-2.85) {31};
\node[sel-scout] at (4.95,-2.85) {24};
\node[sel-scout] at (5.50,-2.85) {34};
\node[sel-scout] at (6.05,-2.85) {38};
\node[sel-sc] at (6.60,-2.85) {44};
\node[sel-sc] at (7.15,-2.85) {86};
\node[sel-scmed] at (7.70,-2.85) {88};
\node[sel-sc] at (8.25,-2.85) {91};
\node[sel-sc] at (8.80,-2.85) {98};
\node[sel-sc] at (9.35,-2.85) {48};
\node[sel-sc] at (9.90,-2.85) {66};
\node[sel-scmed] at (10.45,-2.85) {70};
\node[sel-sc] at (11.00,-2.85) {80};
\node[sel-sc] at (11.55,-2.85) {87};
\node[sel-sc] at (12.10,-2.85) {50};
\node[sel-scpiv] at (12.65,-2.85) {71};
\node[sel-sc] at (13.20,-2.85) {85};
\node[sel-sc] at (13.75,-2.85) {89};
\draw[very thick, black!80] (9.075,-3.21) -- (9.075,-2.49);
\draw[very thick, black!80] (11.825,-3.21) -- (11.825,-2.49);
\node[anchor=east, font=\small] at (0.25,-3.80) {對 71 分割};
\node[sel-scout] at (0.55,-3.80) {13};
\node[sel-scout] at (1.10,-3.80) {30};
\node[sel-scout] at (1.65,-3.80) {35};
\node[sel-scout] at (2.20,-3.80) {15};
\node[sel-scout] at (2.75,-3.80) {33};
\node[sel-scout] at (3.30,-3.80) {16};
\node[sel-scout] at (3.85,-3.80) {27};
\node[sel-scout] at (4.40,-3.80) {31};
\node[sel-scout] at (4.95,-3.80) {24};
\node[sel-scout] at (5.50,-3.80) {34};
\node[sel-scout] at (6.05,-3.80) {38};
\node[sel-sc] at (6.60,-3.80) {44};
\node[sel-sc] at (7.15,-3.80) {48};
\node[sel-sc] at (7.70,-3.80) {66};
\node[sel-sc] at (8.25,-3.80) {70};
\node[sel-sc] at (8.80,-3.80) {50};
\node[sel-scpiv] at (9.35,-3.80) {71};
\node[sel-sc] at (9.90,-3.80) {88};
\node[sel-sc] at (10.45,-3.80) {91};
\node[sel-sc] at (11.00,-3.80) {80};
\node[sel-sc] at (11.55,-3.80) {87};
\node[sel-sc] at (12.10,-3.80) {98};
\node[sel-sc] at (12.65,-3.80) {89};
\node[sel-sc] at (13.20,-3.80) {85};
\node[sel-sc] at (13.75,-3.80) {86};
\node[anchor=east, font=\small] at (0.25,-4.75) {排好 $A[12..16]$};
\node[sel-scout] at (0.55,-4.75) {13};
\node[sel-scout] at (1.10,-4.75) {30};
\node[sel-scout] at (1.65,-4.75) {35};
\node[sel-scout] at (2.20,-4.75) {15};
\node[sel-scout] at (2.75,-4.75) {33};
\node[sel-scout] at (3.30,-4.75) {16};
\node[sel-scout] at (3.85,-4.75) {27};
\node[sel-scout] at (4.40,-4.75) {31};
\node[sel-scout] at (4.95,-4.75) {24};
\node[sel-scout] at (5.50,-4.75) {34};
\node[sel-scout] at (6.05,-4.75) {38};
\node[sel-sc] at (6.60,-4.75) {44};
\node[sel-scpiv] at (7.15,-4.75) {48};
\node[sel-sc] at (7.70,-4.75) {50};
\node[sel-sc] at (8.25,-4.75) {66};
\node[sel-sc] at (8.80,-4.75) {70};
\node[sel-scout] at (9.35,-4.75) {71};
\node[sel-scout] at (9.90,-4.75) {88};
\node[sel-scout] at (10.45,-4.75) {91};
\node[sel-scout] at (11.00,-4.75) {80};
\node[sel-scout] at (11.55,-4.75) {87};
\node[sel-scout] at (12.10,-4.75) {98};
\node[sel-scout] at (12.65,-4.75) {89};
\node[sel-scout] at (13.20,-4.75) {85};
\node[sel-scout] at (13.75,-4.75) {86};
\foreach \k in {1,...,25} { \node[idxlbl] at (\k*0.55, -5.35) {\k}; }
\end{tikzpicture}
\diagcaption{\proc{Select} 找 $25$ 個數的第 $13$ 小。 第一層 $x = 38$，往右；第二層
$A[12..25]$ 的 $x = 71$，往左；第三層只剩 $5$ 個數，直接排序。}
\label{fig:sel-mom-trace}
\end{center}

第一層分成 $5$ 組，組中位數 $M = \langle 80, 35, 38, 31, 50 \rangle$；對 $M$ 找第
$\lceil 5/2 \rceil = 3$ 小，$M$ 只有 $5$ 個數，直接排序得到 $x = 38$。 用 $38$ 分割後
$q = 11$，$k = 11 < 13$，往右邊 $A[12..25]$ 找第 $13 - 11 = 2$ 小。 第二層的 $14$ 個數分
成 $A[12..16]$、$A[17..21]$、$A[22..25]$ 三組，最後一組 $50, 71, 85, 89$ 取下中位數
$71$；$M = \langle 88, 70, 71 \rangle$ 的第 $2$ 小是 $x = 71$。 用 $71$ 分割後
$q = 17$，$k = 6 > 2$，往左邊 $A[12..16]$ 找第 $2$ 小。 第三層只剩 $5$ 個數，直接排成
$44, 48, 50, 66, 70$，答案是 $A[13] = 48$。

驗算：$25$ 個數由小到大的前 $13$ 個是 $13, 15, 16, 24, 27, 30, 31, 33, 34, 35, 38, 44,
48$，第 $13$ 個是 $48$。 第一層丟掉了 $11$ 個，第二層從 $14$ 個留下 $5$ 個；第一層往
右的那一邊有 $14$ 個，不超過上一小節算的 $16$。

\subsection{C 程式}

\texttt{select\_mom} 照虛擬碼翻譯，$M$ 用 \texttt{malloc} 配置，找完 $x$ 就
\texttt{free}；\texttt{partition\_around} 把找位置、交換和 Lomuto 分割寫在一起。 程式
印出第~\ref{sec:sel-mom-trace}~節那 $25$ 個數每一層的追蹤，\texttt{d=} 後面是遞迴到第幾
層的 $M$，\texttt{d=0} 是原陣列。

\begin{example}[label=ex:sel-mom]{c}{(mom\_select.c) 中位數的中位數}
#include <stdio.h>
#include <stdlib.h>

void swap(int *x, int *y)
{
    int t = *x;
    *x = *y;
    *y = t;
}

void insertion_sort(int *a, int l, int h)       // 排好 a[l..h]
{
    for (int i = l + 1; i <= h; i++) {
        int key = a[i], j = i - 1;
        while (j >= l && a[j] > key) {
            a[j + 1] = a[j];
            j--;
        }
        a[j + 1] = key;
    }
}

// 把 x 換到 a[r]，再做 Lomuto 分割，回傳 x 最後的位置。
int partition_around(int *a, int p, int r, int x)
{
    int j = p;
    while (a[j] != x)
        j++;
    swap(&a[j], &a[r]);
    int i = p - 1;
    for (j = p; j < r; j++)
        if (a[j] <= x)
            swap(&a[++i], &a[j]);
    swap(&a[i + 1], &a[r]);
    return i + 1;
}

// 回傳 a[p..r] 裡第 i 小的數，並印出每一層的追蹤。
// d 是遞迴到「中位數陣列」的層數，原陣列是 0。
int select_mom(int *a, int p, int r, int i, int d)
{
    int n = r - p + 1;
    if (n <= 5) {
        insertion_sort(a, p, r);
        printf("d=%d n=%d i=%d base -> %d\n",
               d, n, i, a[p + i - 1]);
        return a[p + i - 1];
    }
    int g = (n + 4) / 5;
    int *m = malloc(g * sizeof *m);
    if (m == NULL)
        exit(1);
    for (int j = 0; j < g; j++) {
        int l = p + 5 * j, h = l + 4 < r ? l + 4 : r;
        insertion_sort(a, l, h);
        m[j] = a[(l + h) / 2];      // 這一組的（下）中位數
    }
    printf("d=%d n=%d i=%d medians", d, n, i);
    for (int j = 0; j < g; j++)
        printf(" %d", m[j]);
    printf("\n");
    int x = select_mom(m, 0, g - 1, (g + 1) / 2, d + 1);
    free(m);
    int k = partition_around(a, p, r, x) - p + 1;
    printf("d=%d x=%d k=%d %s\n", d, x, k,
           i == k ? "found" : (i < k ? "go left" : "go right"));
    if (i == k)
        return x;
    else if (i < k)
        return select_mom(a, p, p + k - 2, i, d);
    else
        return select_mom(a, p + k, r, i - k, d);
}

int main(void)
{
    int a[] = {48, 89, 80, 91, 44, 13, 35, 30, 85, 66, 98, 33, 38,
               88, 15, 70, 87, 31, 16, 27, 24, 50, 86, 71, 34};
    printf("answer %d\n", select_mom(a, 0, 24, 13, 0));
    return 0;
}
\end{example}
\iosample{}{d=0 n=25 i=13 medians 80 35 38 31 50 \\ d=1 n=5 i=3 base -> 38 \\ d=0 x=38 k=11 go right \\ d=0 n=14 i=2 medians 88 70 71 \\ d=1 n=3 i=2 base -> 71 \\ d=0 x=71 k=6 go left \\ d=0 n=5 i=2 base -> 48 \\ answer 48}

追蹤的七行和圖~\ref{fig:sel-mom-trace} 一一對應。 這個演算法每一層都要排好每一組、再遞迴找
$x$，比較次數的常數比範例~\ref{ex:sel-rselect} 的大約 $3.3n$ 大得多。 所以實務上多半用隨機化
選擇，或者先用隨機化選擇、發現遞迴太深時才換成中位數的中位數（Musser 在 1997 年提出的
introselect）。 中位數的中位數的價值在於它證明了
「最壞情況線性」做得到，而它的分析方法（兩個遞迴的比例加起來小於 $1$）在其他問題也
用得上。

\section{本章常見錯誤}
\label{sec:sel-pitfalls}

\begin{center}
\begin{tabular}{p{0.25\linewidth}p{0.50\linewidth}p{0.15\linewidth}}
\toprule
錯誤 & 例子 & 見 \\
\midrule
往右走時名次沒有減 & 遞迴寫成 \proc{Randomized-Select}$(A, q+1, r, i)$，答案錯，或
範圍外存取 & 第~\ref{sec:sel-random}~節 \\
名次和索引差一 & 題目的 $k$ 從 $1$ 開始，C 陣列從 $0$ 開始，直接拿 $k$ 當索引 &
第~\ref{sec:sel-trace}~節 \\
兩邊都遞迴 & 寫成快速排序，時間變成 $\Theta(n \lg n)$ & 第~\ref{sec:sel-random}~節 \\
固定取最後一個當基準 & 已經排好的輸入變成 $\Theta(n^2)$ & 第~\ref{sec:sel-random}~節 \\
忽略重複值 & 全部相等時 Lomuto 分割每次只少一個數，隨機也沒用 &
第~\ref{sec:sel-random}~節 \\
以為 $x$ 是中位數 & 中位數的中位數只保證兩邊各有至少約 $3n/10$ 個數 &
第~\ref{sec:sel-mom}~節 \\
3 個一組也是線性 & 遞迴式 $T(n/3) + T(2n/3) + an$ 的解是 $\Theta(n \lg n)$ &
第~\ref{sec:sel-mom}~節 \\
決賽輸家當第二大 & 第二大可能第一輪就遇到冠軍 & 第~\ref{sec:sel-minmax}~節 \\
最大最小的次數算錯 & 寫成 $3n/2$，沒有分奇偶，或忘了偶數時開頭那一次 &
第~\ref{sec:sel-minmax}~節 \\
\bottomrule
\end{tabular}
\end{center}

\section{觀念題}
\label{sec:sel-exercises}

題型和研究所考試、期中期末考的筆試相近，有選擇、計算、追蹤和證明。 解答在本章最後。

\begin{enumerate}
\item \label{sel-q-tf}（選擇，可能不只一個答案）下列哪些敘述正確？
\begin{enumerate}
\item \proc{Randomized-Select} 的最壞執行時間是 $\Theta(n^2)$。
\item \proc{Select} 第 13 行找到的 $x$ 一定是 $A[p..r]$ 的中位數。
\item 把 \proc{Select} 改成 $3$ 個一組，第~\ref{sec:sel-mom}~節的分析仍然能證明最壞
$O(n)$。
\item 任何只用比較找最大值的演算法，至少要比 $n - 1$ 次。
\item $n = 10$ 時，同時找最大和最小，最壞情況至少要比 $13$ 次。
\end{enumerate}

\item \label{sel-q-tour}（追蹤）
\begin{enumerate}
\item 用錦標賽找 $\langle 93, 41, 67, 28, 85, 52, 76, 39, 60 \rangle$ 的第二大。 相鄰
兩個配對，人數是奇數時最後一個輪空。 畫出比賽樹，寫出直接輸給冠軍的數，以及總比較
次數。
\item 同樣的數改成 $\langle 41, 67, 28, 85, 52, 76, 39, 60, 93 \rangle$，再做一次。 比
較次數為什麼和定理~\ref{thm:sel-second} 不一樣？ 這和下界矛盾嗎？
\end{enumerate}

\item \label{sel-q-trace}（追蹤）用基準固定取最後一個的選擇（把
\proc{Randomized-Partition} 換成 \proc{Lomuto-Partition}），在
$\langle 35, 81, 22, 64, 17, 90, 48, 73, 56 \rangle$ 裡找第 $3$ 小。 每次分割之後寫出
陣列、$p$、$r$、$q$、$k$ 和新的 $i$，並算出總比較次數。

\item \label{sel-q-worst}（計算）
\begin{enumerate}
\item 基準固定取最後一個，在已經由小到大排好的 $n$ 個數裡找最小值，比較幾次？
\item 改用 \proc{Randomized-Select} 找最小值，比較次數和 (a) 一樣多的機率是多少？
\end{enumerate}

\item \label{sel-q-n3}（計算）$n = 3$、$i = 2$，三個數互不相同。 列出
\proc{Randomized-Select} 所有可能的基準選擇，算出期望比較次數，再用
定理~\ref{thm:sel-rselect-cmp} 的公式驗算。

\item \label{sel-q-min}（證明）用定理~\ref{thm:sel-rselect-cmp} 的公式證明：找最小值
（$i = 1$）的期望比較次數是 $2n - 2H_n$，$H_n = 1 + \frac{1}{2} + \cdots + \frac{1}{n}$。
和範例~\ref{ex:sel-rselect} 印出的 $1987.0$ 對照。

\item \label{sel-q-correct}（證明）假設 $A[p..r]$ 的數互不相同，$1 \le i \le r - p + 1$。
對 $n = r - p + 1$ 用強歸納法，證明 \proc{Randomized-Select}$(A, p, r, i)$ 回傳
$A[p..r]$ 裡第 $i$ 小的數。

\item \label{sel-q-dup}（證明）$A[1..n]$ 的 $n$ 個數全部相等，$i = \lfloor (n+1)/2 \rfloor$。
\begin{enumerate}
\item 證明不管 \proc{Random} 傳回什麼，\proc{Randomized-Select} 都要比較
$\Theta(n^2)$ 次。
\item 改寫分割，讓它傳回兩個位置 $q_1 \le q_2$，使 $A[p..q_1 - 1] < x$、
$A[q_1..q_2] = x$、$A[q_2 + 1..r] > x$，再寫出對應的選擇程序（虛擬碼）。 說明為什麼
全部相等時只要比 $O(n)$ 次。
\end{enumerate}

\item \label{sel-q-mom15}（追蹤）用 \proc{Select} 找
$\langle 42, 17, 93, 58, 26, 71, 35, 84, 10, 63, 49, 77, 22, 96, 31 \rangle$ 的中位數。
寫出各組排好的結果、$M$、$x$、分割後的陣列和 $k$。 圖~\ref{fig:sel-mom-cols} 那種
「保證 $\ge x$」的數有幾個？

\item \label{sel-q-quarter}（證明）把 \proc{Select} 第 13 行改成取 $M$ 的第
$\lceil g/4 \rceil$ 小，其他不變。 仿照引理~\ref{lem:sel-mom-split} 估計遞迴範圍的上
界，寫出遞迴式，說明為什麼定理~\ref{thm:sel-mom} 的代換法證不過。

\end{enumerate}

\section{歷屆試題}
\label{sec:sel-pastexam}

以下是研究所入學考試的原題，題目照原卷（語言、記號、配分照錄）。 詳解在本章最後。

\begin{pastexam}[label={px:sel-1}]{中正資工所 113 學年度 第 20 題}
(2\%) Determine if the following statement is true or false: "Find the median of $n$ elements takes $\Omega(nlogn)$ time for any algorithms working in the comparison-based model".
\end{pastexam}

\begin{pastexam}[label={px:sel-2}]{臺大資工所 113 學年度 第 1 題}
（複選題，每個選項單獨計分）

(5 points) Consider the following variation of the Selection Problem: Given a sequence $S=\langle a_1,a_2,\ldots,a_n\rangle$ of $n$ distinct integers and two positive integers $x$ and $y$, where each number in $S$ falls between 1 and $y$, the goal is to report the smallest $x$ numbers in $S$ in ascending order. This scenario assumes the conventional single-processor, random access machine (RAM) model, with a machine word size of at least $\lceil\log y\rceil$ bits. To tackle this problem, five algorithms have been proposed. Please select the correct description(s) below.
\begin{enumerate}[label=\Alph*.]
\item Algorithm A: The algorithm initializes an empty max-heap and sequentially inserts numbers from $S$. During insertion, if the size of the max-heap exceeds $x$, it removes the maximum element from the heap. Finally, the remaining numbers in the heap are sorted by heapsort and reported. If $x\le n/\log n$, Algorithm A runs in $O(n)$ time in the worst case.
\item Algorithm B: The algorithm initially constructs a min-heap and proceeds to extract the minimum element $x$ times. If $x\le n/\log n$, Algorithm B runs in $O(n)$ time in the worst case.
\item Algorithm C: The algorithm begins by performing a base-$n$ radix sort, followed by reporting the smallest $x$ values. If $y\le n^2$, Algorithm C runs in $O(n)$ time in the worst case.
\item Algorithm D: The algorithm performs a counting sort by allocating an array $C$ of length $y$, traversing through $S$ and incrementing $C[a_i]$ for each $a_i$. It finally reports the smallest $x$ values by scanning through the counter array. If $y\le n^2$, Algorithm D runs in $O(n)$ time in the worst case.
\item Algorithm E: The algorithm first applies quickselect, i.e., Randomized-Select, to determine the $x$-th smallest number in $S$. Subsequently, it traverses $S$ for extracting all numbers smaller than or equal to the $x$-th smallest number. Finally, these $x$ numbers are sorted by quicksort and reported. The time complexity of Algorithm E is $O(n+x\log x)$ in expectation.
\end{enumerate}
\end{pastexam}

\begin{pastexam}[label={px:sel-3}]{中央資工所 115 學年度 第 10 題}
（多選題，每一個選項單獨計分，每題 5 分）

Consider the prune-and-search algorithm using the median of medians for finding the $k$-th smallest element in an unsorted array $A$ of $n$ distinct elements, where $1\le k\le n$, as described below:

\begin{quote}
Algorithm Select($A$, $k$):\\
Input: An array $A$ of $n$ elements and an integer $k$, $1\le k\le |A|=n$\\
Output: The $k$-th smallest element in $A$\\
Step 1: If $|A|\le 5$, then sort array $A$ (with any method, e.g., insertion sort) and return the element with rank $k$ (i.e., the $k$-th element in the sorted array).\\
Step 2: Divide the array $A$ into groups of five elements each (ignore the last group if it has fewer than five elements).\\
Step 3: Find the median of each group.\\
Step 4: Recursively compute the median of the found medians, denoted by $m$.\\
Step 5: Partition the original array into the following three sets:\\
$L=\{x\in A\mid x<m\}$\\
$E=\{m\}$\\
$R=\{x\in A\mid x>m\}$\\
Step 6: If $k=|L|+1$ then return $m$\\
\hspace*{3.2em}Else if $k\le|L|$ then return Select$(L,k)$\\
\hspace*{3.2em}Else return Select$(R,k-|L|-1)$.
\end{quote}

Select one or more correct statements.
\begin{enumerate}[label=\Alph*.]
\item After grouping elements into blocks of five and choosing the median of medians as the pivot denoted $m$, at least $3n/10$ elements are guaranteed to be less than or equal to the pivot.
\item The recurrence $\mathrm{T}(n)=\mathrm{T}(n/5)+\mathrm{T}(7n/10)+\mathrm{O}(n)$ implies that the algorithm runs in $\mathrm{O}(n\log n)$ time.
\item Using groups of size 3 instead of 5 still guarantees linear worst-case running time.
\item The correctness of the pruning step relies on the fact that the pivot is an approximate median of the entire array.
\item The median-finding problem is a special case of the $k$-th smallest selection problem, where $k=\lceil n/2\rceil$, and both can be solved in worst-case linear time.
\end{enumerate}
\end{pastexam}

\begin{pastexam}[label={px:sel-4}]{成大資訊聯招 109 學年度 演算法第 3 題}
Show how the time complexity of quicksort can be made to run in $O(n\lg n)$ time. \quad (10\%)
\end{pastexam}

\section{實作練習}
\label{sec:sel-practice}

下面是 Aki-Judge 上和本章相關的題目。 兩題都只交函式（第二題要交標頭檔和 \texttt{.c}
檔），回傳第 $k$ 小的值（$k$ 從 $1$ 開始），不寫 \texttt{main}。

\begin{center}
\begin{tabular}{llc}
\toprule
小節 & 題目 & 難度 \\
\midrule
隨機化選擇 & \akiref{k-th-smallest-element}{第 K 小的元素} & 基本 \\
中位數的中位數 & \akiref{worst-case-linear-k-th-smallest-function}{最壞情況線性時間的第 k 小函式} & 進階 \\
\bottomrule
\end{tabular}
\end{center}

第一題的 $n$ 到 $10^6$。 固定取最後一個當基準，遇到已經排好的輸入會變成
$\Theta(n^2)$；大量重複的值也會讓 Lomuto 分割變慢（第~\ref{sec:sel-random}~節最後一
小節）。 寫的時候用隨機基準，再加上三路分割或第~14~章的 Hoare 分割，程式對任何輸入
都不會退化。 第二題請照
第~\ref{sec:sel-mom}~節寫中位數的中位數；注意題目允許重複值，分割時等於 $x$ 的數不能
都擠到同一邊，否則引理~\ref{lem:sel-mom-split} 的保證就不成立。 交之前先用
範例~\ref{ex:sel-mom} 的資料比對每一層的結果。

\section{讀完你應該會}
\label{sec:sel-checklist}

\begin{itemize}
\item 能寫出成對處理的最大最小演算法，證明它剛好比 $\lceil 3n/2 \rceil - 2$ 次，並說
出對手論證為什麼不能更少。
\item 能畫出錦標賽的比賽樹，找出第二大，算出 $n + \lceil \lg n \rceil - 2$ 次的由
來。
\item 能在給定的分割方法下，逐次寫出選擇演算法每一次分割後的陣列、$q$、$k$ 和新的
名次。
\item 能用 $U(n)$ 的遞迴式和代換法證明隨機化選擇的期望時間是線性，並用指示變數寫出
期望比較次數。
\item 能說明中位數的中位數為什麼保證兩邊都不超過約 $7n/10$，寫出遞迴式並用代換法證
明線性；能說明 $3$ 個一組為什麼不行。
\end{itemize}

\textbf{延伸閱讀}（書末會整理完整的參考資料）

\begin{itemize}
\item Thomas H. Cormen 等，\textit{Introduction to Algorithms}，第 4 版，第 9 章：本章
三節的對應內容。 它證明隨機化選擇的方法和本章不同（數「有幫助的分割」），中位數的
中位數也用了另一種分組方式，可以對照著看。
\item Donald E. Knuth，\textit{The Art of Computer Programming, Vol.~3}，第 2 版，
5.3.3 節：用最少的比較次數找第 $i$ 小，包括第二大的下界的完整證明。
\item M. Blum、R. W. Floyd、V. Pratt、R. L. Rivest、R. E. Tarjan，``Time Bounds for
Selection''，\textit{Journal of Computer and System Sciences} 7(4)，1973：中位數的中
位數的原始論文，也給了選擇問題比較次數的下界。
\item K. Chen、A. Dumitrescu，``Select with Groups of 3 or 4''，WADS 2015：改變挑基準
的方式之後，$3$ 個或 $4$ 個一組也能做到最壞線性時間。
\end{itemize}

\addcontentsline{toc}{section}{觀念題解答}
\section*{觀念題解答}

\begin{enumerate}
\item (a)、(d)、(e) 正確。 (a)：每次都挑到最壞的基準時是 $n(n-1)/2$ 次比較。 (b)：$x$
只保證兩邊各至少約 $3n/10$ 個數，第~\ref{sec:sel-mom-trace}~節的 $x = 38$ 是第 $11$
小，不是第 $13$ 小。 (c)：遞迴式變成 $T(n/3) + T(2n/3) + an$，比例加起來是 $1$，只能
證出 $O(n \lg n)$。 (d)：除了最大值，每個數都要輸一場。 (e)：$\lceil 30/2 \rceil - 2 = 13$。

\item
\begin{enumerate}
\item 第一輪：$93 > 41$、$67 > 28$、$85 > 52$、$76 > 39$，$60$ 輪空。 第二輪：
$93 > 67$、$85 > 76$，$60$ 輪空。 第三輪：$93 > 85$，$60$ 輪空。 第四輪：$93 > 60$。
找冠軍比了 $4 + 2 + 1 + 1 = 8$ 次。 直接輸給 $93$ 的是 $41, 67, 85, 60$，再比 $3$ 次
得到第二大 $85$。 共 $11 = 9 + \lceil \lg 9 \rceil - 2$ 次。
\item 第一輪：$67$、$85$、$76$、$60$ 晉級，$93$ 輪空；第二輪：$85$、$76$ 晉級，$93$
輪空；第三輪：$85$ 晉級，$93$ 輪空；第四輪：$93 > 85$。 找冠軍一樣比 $8$ 次，但冠軍只
比了一場，直接輸給它的只有 $85$，不用再比，共 $8$ 次。 定理~\ref{thm:sel-second} 說
的是最壞情況的次數；這個輸入剛好讓冠軍一路輪空，比較少。 下界也是對最壞情況說的：
任何演算法都\textbf{存在}某個輸入要比 $11$ 次，不是每個輸入都要，所以不矛盾。
\end{enumerate}

\item
\begin{center}
\begin{tabular}{cccccl}
\toprule
$p..r$ & $i$ & 基準 & 分割後的陣列 & $q$, $k$ & 下一步 \\
\midrule
$1..9$ & 3 & 56 & 35 22 17 48 56 90 64 73 81 & 5, 5 & 往左 \\
$1..4$ & 3 & 48 & 35 22 17 48 56 90 64 73 81 & 4, 4 & 往左 \\
$1..3$ & 3 & 17 & 17 22 35 48 56 90 64 73 81 & 1, 1 & 往右，$i \gets 2$ \\
$2..3$ & 2 & 35 & 17 22 35 48 56 90 64 73 81 & 3, 2 & 回傳 $35$ \\
\bottomrule
\end{tabular}
\end{center}
比較次數 $8 + 3 + 2 + 1 = 14$。

\item (a) 每次的基準都是範圍裡最大的數，$k = n$，往左邊少一個數的範圍找。 比較次數
$(n-1) + (n-2) + \cdots + 1 = n(n-1)/2$。 (b) 要比這麼多次，範圍是 $m$ 個數時基準必
須是其中最大的（任何其他基準都會一次丟掉不只一個數），機率 $1/m$。 $m = n, n-1,
\ldots, 2$ 每一步獨立，機率是 $\frac{1}{n} \cdot \frac{1}{n-1} \cdots \frac{1}{2} =
\frac{1}{n!}$。

\item 第一次分割一定比 $2$ 次。 基準是第 $2$ 小（機率 $1/3$）：直接找到，共 $2$ 次。
基準是第 $1$ 小（機率 $1/3$）：往右邊兩個數找第 $1$ 小，再分割一次比 $1$ 次，不管基
準是哪個，下一步不是找到就是只剩一個數，共 $3$ 次。 基準是第 $3$ 小：對稱，$3$ 次。
期望值 $(2 + 3 + 3)/3 = 8/3$。 公式：$(a, b) = (1, 2)$ 的分母是 $2 - 1 + 1 = 2$，
$(1, 3)$ 是 $3$，$(2, 3)$ 是 $2$，和是 $1 + \frac{2}{3} + 1 = \frac{8}{3}$，一樣。

\item $i = 1$ 時 $\min(a, 1) = 1$、$\max(b, 1) = b$，每一對的機率是 $2/b$。 固定 $b$，
$a$ 有 $b - 1$ 種：
\[
E_n(1) = \sum_{b=2}^{n} \frac{2(b-1)}{b} = 2(n - 1) - 2\sum_{b=2}^{n} \frac{1}{b}
= 2(n - 1) - 2(H_n - 1) = 2n - 2H_n.
\]
$n = 1001$：$H_{1001} \approx \ln 1001 + 0.5772 \approx 7.486$，$2002 - 14.97 \approx
1987.0$，和程式印的公式值一樣。

\item $n = 1$：$p = r$，$i$ 只能是 $1$，回傳唯一的數，正確。 $n \ge 2$：假設比 $n$ 小
的範圍都正確。 \proc{Lomuto-Partition} 結束後 $A[p..q-1]$ 都 $< A[q]$、$A[q+1..r]$ 都
$> A[q]$（數互不相同），而且 $A[p..r]$ 仍是原來那些數，所以 $A[q]$ 是第 $k = q - p + 1$
小。 $i = k$ 時回傳 $A[q]$，正確。 $i < k$ 時，第 $i$ 小的數比 $A[q]$ 小，在
$A[p..q-1]$ 裡，而且這一段的數正是全部裡最小的 $k - 1$ 個，所以它在這一段裡也是第 $i$
小；範圍有 $k - 1 \ge i$ 個數、小於 $n$，由歸納假設遞迴回傳正確。 $i > k$ 時，答案在
$A[q+1..r]$，比它小的數裡有 $k$ 個在左邊（含基準），所以它在右邊是第 $i - k$ 小，範圍
有 $n - k \ge i - k \ge 1$ 個數，同樣由歸納假設正確。

\item (a) 所有數相等，隨機交換不改變陣列。 Lomuto 分割裡每個 $A[j] \le x$ 都成立，結
束時 $q = r$、$k = n$。 $i < k$ 時往左邊 $n - 1$ 個數遞迴，名次不變。 範圍從 $n$ 一路
減一，直到範圍大小等於 $i$ 時 $k = i$ 才停，比較次數
$\sum_{m=i}^{n} (m - 1) \ge (n - i + 1)(i - 1)$。 $i = \lfloor (n+1)/2 \rfloor$ 時這是
$\Theta(n^2)$。

(b) 三路分割的一種寫法：用 $\mathit{lt}$、$j$、$\mathit{gt}$ 三個指標，維持
$A[p..\mathit{lt}-1] < x$、$A[\mathit{lt}..j-1] = x$、$A[\mathit{gt}+1..r] > x$，$j$
從左往右掃：$A[j] < x$ 就和 $A[\mathit{lt}]$ 交換、兩者都加一；$A[j] > x$ 就和
$A[\mathit{gt}]$ 交換、$\mathit{gt}$ 減一（$j$ 不動）；相等就只把 $j$ 加一。 掃完傳
回 $q_1 = \mathit{lt}$、$q_2 = \mathit{gt}$。 選擇程序和 \proc{Randomized-Select} 一樣，
只是令 $k_1 = q_1 - p + 1$、$k_2 = q_2 - p + 1$：$i < k_1$ 往左找 $A[p..q_1 - 1]$ 的第
$i$ 小，$i > k_2$ 往右找 $A[q_2 + 1..r]$ 的第 $i - k_2$ 小，否則回傳 $x$。
全部相等時第一次分割就得到 $q_1 = p$、$q_2 = r$，任何 $i$ 都落在中間那段，分割一次就
回傳，比較 $O(n)$ 次。

\item 三組排好：$\langle 17, 26, 42, 58, 93 \rangle$、$\langle 10, 35, 63, 71, 84
\rangle$、$\langle 22, 31, 49, 77, 96 \rangle$，$M = \langle 42, 63, 49 \rangle$，$M$
的第 $2$ 小是 $x = 49$。 用 $49$ 分割後陣列是 $\langle 17, 26, 42, 10, 35, 22, 31, 49,$
$71, 84, 58, 93, 96, 77, 63 \rangle$，$q = 8$，$k = 8 = i$，第一層就找到中位數 $49$。
保證 $\ge 49$ 的：中位數 $\ge 49$ 的兩組（中位數 $63$、$49$）各 $3$ 個，即
$63, 71, 84, 49, 77, 96$，共 $6$ 個；保證 $\le 49$ 的也是 $6$ 個（$17, 26, 42, 22, 31,
49$）。

\item 中位數 $\ge x$ 的組有 $g - \lceil g/4 \rceil + 1$ 組，大約 $3g/4$，保證 $\ge x$
的數大約 $3 \cdot 3g/4 = 9n/20$，左邊最多大約 $11n/20$。 中位數 $\le x$ 的組只有
$\lceil g/4 \rceil$ 組，保證 $\le x$ 的大約 $3g/4 = 3n/20$，右邊最多大約
$17n/20 + O(1)$。 遞迴式是 $T(n) \le T(\lceil n/5 \rceil) + T(17n/20 + O(1)) + an$。
代換法：$c\frac{n}{5} + c\frac{17n}{20} + an = \frac{21}{20}cn + an > cn$，不管 $c$ 多
大都證不出 $T(n) \le cn$。 原因是 $\frac{1}{5} + \frac{17}{20} = \frac{21}{20} > 1$，遞
迴樹每往下一層，總工作量反而變大。 （用第~2~章第~\ref{sec:rec-gf}~節的 Akra--Bazzi 方法，
$(1/5)^p + (17/20)^p = 1$ 的解 $p \approx 1.12$，這個上界只能給出
$O(n^{1.12})$。）

\end{enumerate}

\addcontentsline{toc}{section}{歷屆試題詳解}
\section*{歷屆試題詳解}

\pastsol{px:sel-1}{中正資工所 113 學年度 第 20 題}
答案：False。

\textbf{上界。} 第~\ref{sec:sel-mom}~節的 \proc{Select} 只透過比較元素的大小得知資訊（各組的插入
排序、找 $x$ 的位置、Lomuto 分割都只比較元素），是比較模型裡的演算法；由定理~\ref{thm:sel-mom}，
它找第 $\lfloor (n+1)/2 \rfloor$ 小的最壞時間是 $O(n)$。 既然有比較模型裡的 $O(n)$ 演算法，
「任何」演算法都要 $\Omega(n \lg n)$ 就不對。

\textbf{為什麼排序的下界搬不過來。} 排序的 $\Omega(n \lg n)$ 來自「答案有 $n!$ 種」（第~15~章的
決策樹）。 找中位數只要回答「是哪一個元素」，答案只有 $n$ 種，決策樹只能給出 $\lceil \lg n
\rceil$ 的下界。

\textbf{真正的下界是線性。} 設演算法回答 $m$，而另一個元素 $y$ 一次都沒有參加比較。 $y < m$
時把 $y$ 改成比所有數都大，$y > m$ 時改成比所有數都小；所有比較的結果都不變，演算法仍然回答
$m$，但比 $m$ 小的數的個數變了，$m$ 不再是中位數。 所以除了 $m$，每個元素都至少要參加一次
比較，一次比較只碰兩個元素，至少要 $(n-1)/2$ 次。 找中位數在比較模型裡是 $\Theta(n)$。

\pastsol{px:sel-2}{臺大資工所 113 學年度 第 1 題}
答案：B、C、E。 以下 $\log$ 的底數不影響結論。
\begin{itemize}
\item A 錯誤。 堆積最多 $x + 1$ 個數，每次插入、刪除最大值是 $O(\log x)$，總共
$O(n \log x)$，最後的堆積排序 $O(x \log x)$。 問題在 $O(n \log x)$ 不一定是 $O(n)$。 取
$x = \lfloor n/\log n \rfloor$、$S$ 由小到大排好：堆積裝滿 $x$ 個之後，每個新來的數都比堆積
裡全部的數大，插入時要從新的葉節點一路上浮到根，交換 $\lfloor \log(x+1) \rfloor$ 次。 後面
$n - x \ge n/2$ 個數各花這麼多，共至少 $\frac{n}{2}\lfloor \log(x+1) \rfloor$。 $n \ge 16$ 時
$x + 1 \ge n/\log n \ge \sqrt{n}$，$\log(x+1) \ge \frac{1}{2}\log n$，所以是 $\Omega(n \log n)$，
不是 $O(n)$。
\item B 正確。 由下而上建最小堆積 $O(n)$（第~8~章第~\ref{sec:hp-build}~節），取出最小值 $x$ 次，
每次 $O(\log n)$，依取出的順序就是由小到大。 合計 $O(n + x \log n) \le
O(n + \frac{n}{\log n} \cdot \log n) = O(n)$。
\item C 正確。 每個數 $a_i \le y \le n^2$，以 $n$ 為基底最多 $3$ 位（$n^2$ 本身是 $100$）。
LSD 基數排序做 $3$ 趟，每趟是 $n$ 個數、$n$ 種位值的計數排序，$\Theta(n + n)$
（第~\ref{sec:lin-lsd}~節「一位要切多大」的最後一段）。 字組至少 $\lceil \log y \rceil$ 位元，
$a_i$ 放得進一個字組，取出某一位（除以 $n^t$ 再取餘數）是常數時間。 再輸出前 $x$ 個，
$O(x)$。 合計 $O(n)$。
\item D 錯誤。 光是配置並掃過長度 $y$ 的陣列 $C$ 就要 $\Theta(y)$；條件只說 $y \le n^2$，
$y = n^2$ 時是 $\Theta(n^2)$，不是 $O(n)$。 （$y = O(n)$ 時才是線性，見第~\ref{sec:lin-counting}~節的成本分析。）
\item E 正確。 \proc{Randomized-Select} 的期望時間是 $O(n)$（定理~\ref{thm:sel-rselect}）；
掃一次 $S$ 取出不大於第 $x$ 小的數是 $O(n)$，數兩兩不同，剛好取出 $x$ 個；（隨機化的）快速
排序排 $x$ 個數，期望 $O(x \log x)$（第~14~章定理~\ref{thm:so-quick-avg}）。 由期望值的線性，
總期望時間 $O(n + x \log x)$。
\end{itemize}

\pastsol{px:sel-3}{中央資工所 115 學年度 第 10 題}
答案：A、E。

先說明解讀。 (A) 和 (B) 照字面嚴格讀，結論會和出題的本意相反：(A) 的 $3n/10$ 要扣掉一個
常數才嚴格成立，(B) 的 $O(n \log n)$ 字面上也是對的上界。 這裡照出題的本意判斷（(A) 是課
本常見的近似說法；(B) 問的是這個遞迴式給出的界），下面把嚴格的情況也寫出來。
\begin{itemize}
\item A（照本意）正確。 Step 2 只用滿的組，共 $g = \lfloor n/5 \rfloor$ 組。 中位數不大於 $m$ 的組
至少 $\lceil g/2 \rceil$ 組（包括 $m$ 自己那組），每組的中位數和比它小的兩個數，共 $3$ 個都
$\le m$，所以至少 $3\lceil g/2 \rceil \approx 3n/10$ 個數 $\le m$。 嚴格地說，這個數可能比
$3n/10$ 少一點：$n = 11$，$A = \langle 1, 2, 3, 20, 21, 10, 11, 12, 22, 23, 30 \rangle$，最後的
$30$ 被忽略，兩組的中位數是 $3$、$12$，取下中位數 $m = 3$，$\le m$ 的只有 $3$ 個，比
$3n/10 = 3.3$ 少。 一般來說 $3\lceil g/2 \rceil \ge \frac{3}{2}\lfloor \frac{n}{5} \rfloor
\ge \frac{3(n-4)}{10}$，所以保證的是至少 $\frac{3n}{10} - \frac{6}{5}$ 個。
\item B（照本意）錯誤。 $\frac{1}{5} + \frac{7}{10} = \frac{9}{10} < 1$，這個遞迴式的解是
$O(n)$（定理~\ref{thm:sel-mom} 的代換法，或第~\ref{sec:sel-mom}~節的遞迴樹：每一層的工作量是上
一層的 $9/10$）。 演算法是線性時間，說它「是 $O(n \log n)$」沒有給出它推得出的界。 嚴格地
說，$O(n)$ 的函數也是 $O(n \log n)$，所以字面上這句話不算錯；作答時可以註明這一點。
\item C 錯誤。 $3$ 個一組時，保證不大於 $m$ 的數只有大約 $2 \cdot \frac{g}{2} = \frac{n}{3}$
個，往下遞迴的一邊可能有 $\frac{2n}{3}$ 個，遞迴式是 $T(n) = T(n/3) + T(2n/3) + O(n)$，比例
加起來是 $1$，解是 $\Theta(n \lg n)$（第~\ref{sec:sel-mom}~節「為什麼是 5 個一組」）。 不能保
證線性。
\item D 錯誤。 正確性不需要 $m$ 接近中位數：Step 5 之後 $L$ 的數都比 $m$ 小、$R$ 的都比 $m$
大，$m$ 是第 $|L| + 1$ 小；$k$ 落在哪一段、在那一段裡是第幾小都由此決定，不管 $m$ 是誰都對
（和本章觀念題第~\ref{sel-q-correct}~題的歸納法一樣）。 $m$ 接近中位數只影響\textbf{時間}：它
保證 $L$、$R$ 都不超過大約 $7n/10$。
\item E 正確。 中位數就是第 $\lceil n/2 \rceil$ 小（$n$ 是偶數時是下中位數，和本書的約定相同），
是選擇問題 $k = \lceil n/2 \rceil$ 的特例；\proc{Select} 對任何 $k$ 的最壞時間都是 $O(n)$
（定理~\ref{thm:sel-mom}）。
\end{itemize}

\pastsol{px:sel-4}{成大資訊聯招 109 學年度 演算法第 3 題}
快速排序本來的期望時間就是 $O(n \lg n)$（第~14~章定理~\ref{thm:so-quick-avg}），所以題目要的
是\textbf{最壞情況}也是 $O(n \lg n)$。 最壞情況發生在每次都切得極偏；如果每次都用
\textbf{中位數}當基準，兩邊各剩大約一半，而第~\ref{sec:sel-mom}~節的 \proc{Select} 可以在最壞
$O(n)$ 時間找到中位數。 以下假設鍵互不相同。

\begin{pseudo}{成大 109 學年度演算法第 3 題的解答}
\Procedure{Median-Quicksort}{$A, p, r$}
    \If{$p < r$}
        \State $x \gets$ \Call{Select}{$A, p, r, \lfloor (r - p + 2)/2 \rfloor$} \Comment{$A[p..r]$ 的中位數}
        \State $q \gets$ \Call{Partition-Around}{$A, p, r, x$}
        \State \Call{Median-Quicksort}{$A, p, q - 1$}
        \State \Call{Median-Quicksort}{$A, q + 1, r$}
    \EndIf
\EndProcedure
\end{pseudo}

\proc{Select} 和 \proc{Partition-Around} 都是虛擬碼~\ref{alg:sel-mom} 裡的程序。 範圍有
$n = r - p + 1$ 個數，$\lfloor (r - p + 2)/2 \rfloor = \lfloor (n+1)/2 \rfloor$，就是中位數的名次。

\textbf{正確性。} \proc{Select} 只在 $A[p..r]$ 裡搬動元素，不改變這一段裝的是哪些數。
\proc{Partition-Around} 把 $x$ 換到 $A[r]$ 再做 Lomuto 分割，結束時
$A[p..q-1] < A[q] = x < A[q+1..r]$。 之後和第~14~章的 \proc{Quicksort}（虛擬碼~\ref{alg:so-quick}）
一樣：對 $n$ 歸納，兩邊遞迴的範圍都比 $n$ 小，由歸納假設各自排好，左邊每個數小於 $x$、右邊
每個數大於 $x$，所以整段排好。 快速排序的正確性不依賴基準選誰，換成中位數也一樣。

\textbf{時間。} 設 $T(n)$ 是 $n$ 個數時的最壞時間。 第 3 行由定理~\ref{thm:sel-mom} 是 $O(n)$，
第 4 行找 $x$ 的位置和分割都是 $O(n)$，合計不超過 $cn$。 $x$ 是第 $\lceil n/2 \rceil$ 小，左邊有
$\lceil n/2 \rceil - 1 \le \lfloor n/2 \rfloor$ 個數，右邊有 $n - \lceil n/2 \rceil = \lfloor n/2
\rfloor$ 個數。 所以
\[
T(n) \le 2T(\lfloor n/2 \rfloor) + cn .
\]
用遞迴樹看：第 $d$ 層的每個範圍不超過 $n/2^d$ 個數，同一層的範圍互不重疊，所以每一層的工作量
加起來不超過 $cn$；範圍每往下一層至少減半，$\lfloor \lg n \rfloor + 1$ 層以內就只剩 $0$ 或 $1$
個數。 總時間 $O(n \lg n)$，這也是主定理（第~2~章定理~\ref{thm:rec-master}）$a = b = 2$、
$f(n) = \Theta(n)$ 的情況 2。

\textbf{補充。} 這個做法的常數很大（\proc{Select} 本身就要很多次比較），實務上不會這樣寫。 常見
的做法是隨機選基準，只保證期望 $O(n \lg n)$；要保證最壞情況，C++ 程式庫的 introsort 在遞迴太
深時改用堆積排序（第~14~章第~\ref{sec:so-heap}~節最後）。 作答時寫出上面的演算法和遞迴式就夠了，
可以加一句說明這兩種實務做法。 鍵有重複時，把分割改成三路分割（第~14~章觀念題第 6 題），等於
$x$ 的數全部留在中間，兩邊仍然各不超過一半。
